\documentclass[preprint,12pt]{elsarticle}
\usepackage[T1]{fontenc}
\usepackage{lmodern}
\usepackage{amsmath}
\usepackage{graphicx}
\usepackage{booktabs}
\usepackage[section]{placeins}
\usepackage{longtable}
\usepackage{array}
\usepackage{textcomp}
\usepackage{xurl}
\usepackage[hidelinks]{hyperref}
\newcommand{\degC}{\,\textdegree C}
\renewcommand{\thesection}{S\arabic{section}}
\renewcommand{\thetable}{S\arabic{table}}
\renewcommand{\thefigure}{S\arabic{figure}}
\biboptions{sort&compress}
\setlength{\emergencystretch}{1.5em}
\makeatletter
\def\ps@pprintTitle{\let\@oddhead\@empty\let\@evenhead\@empty\def\@oddfoot{\reset@font\hfil\thepage\hfil}\let\@evenfoot\@oddfoot}
\makeatother

\begin{document}

\begin{frontmatter}
\title{Supplementary material for: Thermal runaway propagation between two lithium-ion cells: resolved and lumped neighbour models against sixteen tests}
\author{Siddharth Satte}
\address{Independent researcher, Pune, India}
\end{frontmatter}

This supplement holds the record of the study (its sixteen steps and the predictions registered before their runs), the numbers behind the figures of the main text test by test, the checks, the inputs of the bounds on what each path can deliver, how the one published history of a propagating neighbour was read, the reduced models' and the held-out test's numbers, and step 11's searches under its kinetic and thermal draws. Every label it uses is defined in Section~\ref{sec:s-key}, so that it can be read without the archive, and the registry quotes the predictions as they were committed, shortened where long (Section~\ref{sec:s-registry}), with the expectations recorded beside them, which are not results of this paper; it lists the predictions whose runs the main text or this supplement reports. References are numbered in this supplement's own list. The archive~\cite{satte2026repo} holds the code, the run records and a log of every decision, for anyone who wants to repeat a result or read a prediction's reasoning in full.

\section{The labels used in the registry}
\label{sec:s-key}

The study's records, and so the predictions in Table~\ref{tab:s-registry}, use short labels. In plain words:
\begin{itemize}
\item Step: one of the study's sixteen stages, each with its own predictions file (Table~\ref{tab:s-steps}); the records call a step a stage, and the registry writes step. Each step had a plan, written before its code and its predictions; the plan's expectation is the outcome the plan expected, a judgement it states without a computation unless the registry quotes one, and the plan's line for a prediction is the line of the plan that states it. The study's expectation is the outcome the predictions file itself recorded before the runs, from earlier steps' records, hand estimates or an energy balance that the file sets out; the predictions file's analysis is that reasoning.
\item Trigger, neighbour: the failing cell, modelled as one temperature, and the cell beside it. A lumped neighbour, lumped cell or lumped node is one temperature; a resolved neighbour or cell is a three-dimensional temperature field.
\item Hottest node, hot spot: the resolved neighbour's hottest point at each time. Mean: the resolved neighbour's heat-capacity-weighted mean temperature. Far face: the resolved neighbour's surface point at mid-length opposite the line of contact, called the thermocouple node or step 9's node in the records, where it stood for the measured thermocouple before the thermocouple's angle was found unreported. The model's four other probes at mid-length are surface nodes at 0\textdegree, on the line of contact, at 0.33\textdegree, the mesh's first surface node off it, and at 88.3 and 94.4\textdegree, its two nodes nearest 90\textdegree, the side probes.
\item Onset: the time at which a node's reaction self-heating first reaches 0.15~K\,min$^{-1}$, the exotherm threshold of the accelerating-rate calorimetry that the kinetic source simulated for its reference cell~\cite{koenig2025uncertainty}. The internal-short band: 255.7 to 302.1\degC, the kinetic source's band of two standard deviations for the temperature of an internal short, read from its figure (Section~\ref{sec:s-steps}).
\item C1, the local criterion: the hot spot reaches 255.7\degC. C2, the lumped criterion, called the whole-cell clock in the records: the lumped neighbour reaches 255.7\degC. C8: the resolved neighbour's mean reaches 255.7\degC. Step 7 also scored five criteria that later steps dropped: C3, the lumped neighbour at 150\degC; C4 and C5, the hot spot and the lumped neighbour at 138.5\degC, the peak of the separator's melting in the posterior-mean kinetics; C6 and C7, C1 and C4 with the aged SEI. A criterion is met when its level is reached within 900~s of the trigger's start.
\item P1, P2 and so on: a step's predictions, numbered in its predictions file; Table~\ref{tab:s-registry} lists them by step and number, so that step 13's P7 is the seventh prediction of step 13.
\item 150\degC{} (steps 5 to 7): the least onset temperature of propagation in Fransson et al.'s tests, read off their Fig.~3a, which they give as the minimum a neighbour must reach for propagation to occur~\cite{fransson2024exploring}; an observation, not their test of propagation, the second spike (Section~2.1 of the main text).
\item $E_\mathrm{b}$: the trigger's body energy; $E_\mathrm{e}=3.105\,E_\mathrm{b}$: the energy its ejecta carry, with 3.105 the median ratio of ejected to body energy in the ten nail tests of this cell in the Battery Failure Databank~\cite{finegan2024databank}; the bracket on that ratio is its smallest and its largest in those ten tests, 2.22 and 4.16, which rescale every fraction by 1.40 and 0.75. $E_\mathrm{fix}=76.74$~kJ: the Databank's median ejected energy, one for every test.
\item $f$: the fraction of $E_\mathrm{e}$ delivered to the neighbour through the half of its can facing the trigger. $f^*$ and $E^*$: a criterion's threshold, the least fraction or energy at which it is met. $f_\mathrm{t}$: the fraction at which a criterion is met at a propagated test's measured time (a time-matching fraction). $f_\mathrm{max}$: the fraction at which the far face peaks at a non-propagated test's measured maximum (a peak-matching fraction).
\item Low, central and high contact: a contact conductance of 200, 2,000 and 3,700~W\,m$^{-2}$\,K$^{-1}$ between the cans.
\item Kinetic draws: the 200 joint draws of the kinetic posterior (Section~2.3 of the main text); the four drawn by rank in step 11 are those ranked 5, 50, 151 and 196, earliest first, by the time of the hot spot's onset in step 3's case (Section~\ref{sec:s-steps}), on the earlier two-cell study's geometry, with the cells 2~mm apart and joined by busbars~\cite{satte2026twocell}, the trigger releasing 33.2~kJ over 30~s with a 500~mW\,K$^{-1}$ busbar and no delivered heat; the step's plan named them the 2.5th, 25th, 75th and 97.5th percentiles by nearest rank, which gives 5 and 50 from below, while 151 and 196 mirror them from above, where that rule would give 150 and 195. Thermal draws: 30 Latin hypercube draws of the roll's published thermal properties, its conductivity across the layers (0.998 to 1.38~W\,m$^{-1}$\,K$^{-1}$), along them (24.9 to 25.8~W\,m$^{-1}$\,K$^{-1}$) and its specific heat (842 to 1110~J\,kg$^{-1}$\,K$^{-1}$)~\cite{oregan2022thermal,chen2021simplified}; a thermophysical standard deviation is a spread over them. The nominal properties, called central in the records: the roll's values in Table~2 of the main text, 1.3 and 25~W\,m$^{-1}$\,K$^{-1}$ across and along the layers and 950~J\,kg$^{-1}$\,K$^{-1}$, inside those ranges though not at their middle. Step 11's four thermal draws are those with the smallest and largest conductivity across the layers (1.008 and 1.375~W\,m$^{-1}$\,K$^{-1}$) and the smallest and largest specific heat (849 and 1109~J\,kg$^{-1}$\,K$^{-1}$), each setting all three properties, which Table~\ref{tab:s-step11} gives; its eight draws are these four and its four kinetic draws. The aged SEI: the SEI's initial amount set to zero in every node and in the lumped node, in place of the kinetic source's fresh value~\cite{koenig2025uncertainty}.
\item Run sets: step 5's 5,248 runs, with the earlier two-cell study's trigger, including its factorial of 4,800 (Section~\ref{sec:s-steps}); step 11's four tests, PT-Exp1, SP-Exp3, PP-Exp3 and ST-Exp2. Step 7's hits and false alarms: the propagated and the non-propagated tests whose hot spot reached 255.7\degC{} with no delivered heat.
\item Test names: M50\_PT\_Exp1 and the like in the records are PT-Exp1 and the like here.
\item Terms in the quoted predictions. The estimates above (step 5): the hand estimates its predictions file set out before the runs. The trigger's 33.2~kJ would raise it 538~K adiabatically, and both cells 269~K if shared; through the contact alone the pair evens out on a time scale of 181~s at the central contact (614 and 132~s at the low and the high); a linear two-node balance puts the lumped neighbour's peak without reaction near 145 to 170\degC, 200 to 250~s after the release; if the line of contact rises a third of the way toward the trigger, the hot spot sits near 190 to 200\degC{} for the first few minutes; and a lumped cell must pass about 220\degC{} before its own reaction can carry it to 255.7\degC. The measured values quoted in step 5: the mean heating rate of propagating neighbours, about 0.4~K\,s$^{-1}$, as Fransson et al.\ state it in their text~\cite{fransson2024exploring}, and the share of the tests that propagated, 9 of 16. Step 3's onset finding (step 5): that over step 3's draws the kinetic posterior spread the hot spot's onset time far more than the thermal properties did, the thermal standard deviation being 0.027 of the kinetic. A check on the mapping (step 8): a prediction that tests only the model's mapping from a test's trigger to its thresholds, that a larger trigger never needs more delivered heat. $E^*$ in step 12, a different quantity from the threshold above: the delivered energy at which the lumped criterion is met at a propagated test's measured time, $f_\mathrm{t}E_\mathrm{fix}$. Both contacts (step 11's P4): the low and the high contact, with the nominal properties. The skin model, the cross-section conduction model, the ignition temperature and the two-body model (step 13): the reduced models of Section~2.5 and Appendix~B of the main text. The cap end (step 14): the end of the cell at its positive cap; the paths that step's plan counts as the cap end's are radiation, gas beside the neighbour and the holder, the first, second and fifth paths of Table~4 of the main text. The matched times (step 14): those of the seven propagated tests whose measured time a run of the lumped neighbour reproduced within 10~s with the 21.4~s pulse; PT-Exp4 and ST-Exp3 end at their thresholds.
\item Labels such as $E^*$ are the records' own and are reused from step to step; each is defined here with its step.
\end{itemize}

\section{The sixteen steps}
\label{sec:s-steps}

The study ran in sixteen steps (Table~\ref{tab:s-steps}). Each step's predictions file was committed and pushed to the public repository before any run of the step's registered list, with the outcome expected for each prediction, and was not changed afterwards; the repository records each file's commit, the time of its push and the hash of each step's run list~\cite{satte2026repo}. Two departures are recorded there. While the eleventh step's code was tested, one full run under a kinetic draw was made before its predictions, PT-Exp4 with 5.04~kJ over 21.39~s under the draw of rank 5, whose lumped neighbour ran away at 191.0~s against 215.6~s under the posterior mean, and its result is stated in them; and while that step's runs went, the scoring was once run on 178 of its 216 records to test a figure's script, and the verdicts it printed were seen, after which, as the repository's history shows, no code, rule or run changed.

\IfFileExists{tables/s1_steps.tex}{\begin{table}[!htbp]
\centering
\scriptsize
\caption{The study's sixteen steps, each with its predictions pushed to the public repository before its runs (the push's date, UTC, 2026): the predictions each registered, and how many held and failed among those Table~\ref{tab:s-registry} lists, whose runs the main text or this supplement reports, and among the others; Section~\ref{sec:s-steps} defines a foreseen failure.}
\label{tab:s-steps}
\setlength{\tabcolsep}{2pt}
\begin{tabular}{rp{0.366\textwidth}lrrrrrrr}
\toprule
 & & & & \multicolumn{4}{c}{Listed in Table~\ref{tab:s-registry}} & \multicolumn{2}{c}{Not listed} \\
\cmidrule(lr){5-8}\cmidrule(lr){9-10}
Step & What it asked & Pushed & Predictions & All & Held & Failed & Foreseen & Held & Failed \\
\midrule
1 & Passive kinetics along the archived histories of the earlier two-cell study: when the chemistry starts at the hot spot, the mean and the lumped node & 10 Oct & 4 & 0 & 0 & 0 & 0 & 4 & 0 \\
2 & The reaction heat fed back, one way, to both neighbours & 10 Oct & 3 & 0 & 0 & 0 & 0 & 2 & 1 \\
3 & The onsets over 200 kinetic draws and 30 thermal draws & 12 Oct & 6 & 0 & 0 & 0 & 0 & 3 & 3 \\
4 & Which submodel of the kinetic posterior sets the lead of the hot spot & 16 Oct & 5 & 0 & 0 & 0 & 0 & 4 & 1 \\
5 & The cans in contact and joined by tabs, against the sixteen measured tests & 21 Oct & 7 & 7 & 1 & 6 & 0 & 0 & 0 \\
6 & Why the model does not propagate: propagation read at 150\degC, the least onset temperature of propagation in Fransson et al.'s tests, and the trigger released over the two measured rise times & 23 Oct & 7 & 0 & 0 & 0 & 0 & 1 & 6 \\
7 & Seven ignition criteria scored test by test, each test with its own trigger & 26 Oct & 8 & 1 & 1 & 0 & 0 & 2 & 5 \\
8 & The delivered fraction of the ejected energy that each test's outcome requires & 26 Oct & 8 & 4 & 3 & 1 & 1 & 2 & 2 \\
9 & The delivered fraction read from each test's time and peak & 29 Oct & 8 & 6 & 4 & 2 & 2 & 1 & 1 \\
10 & The one published history, the thermocouple's angle and held-out scores & 30 Oct & 8 & 3 & 2 & 1 & 0 & 2 & 3 \\
11 & Bands of the earlier steps' main numbers under four kinetic draws chosen by rank and four thermal draws at the extremes & 2 Nov & 6 & 5 & 2 & 3 & 1 & 0 & 1 \\
12 & The trigger held to calorimetry's body energy and one ejected energy for every test & 3 Nov & 4 & 4 & 4 & 0 & 0 & 0 & 0 \\
13 & Reduced models of each threshold: skin conduction and an ignition temperature & 16 Nov & 10 & 8 & 6 & 2 & 0 & 1 & 1 \\
14 & What each path can deliver to the neighbour, bounded from published inputs & 17 Nov & 6 & 6 & 4 & 2 & 2 & 0 & 0 \\
15 & The resolved neighbour's own runaway: a reactive cross-section and three-dimensional searches & 18 Nov & 8 & 8 & 7 & 1 & 0 & 0 & 0 \\
16 & The needs and PT-Exp1's threshold in three dimensions, and a refinement & 24 Nov & 11 & 11 & 9 & 2 & 1 & 0 & 0 \\
\midrule
All & & & 109 & 63 & 43 & 20 & 7 & 22 & 24 \\
\bottomrule
\end{tabular}
\end{table}
}{}

The main text reports results of steps 5 and 7 to 16. Steps 1 to 4 concern when the chemistry starts in an earlier geometry of the same cells, with a gap between them and busbars, and step 6 why the model did not propagate, with propagation read at 150\degC{} (Section~\ref{sec:s-key}) and the trigger's energy released over the two rise times derived from the records, 17.53 and 25.24~s (Section~\ref{sec:s-history}); their runs are not reported, but for step 3's ranking of the kinetic draws, which step 11 used, and six runs of step 6 on which step 12's predictions file drew (Section~\ref{sec:s-registry}). Step 3 ran the earlier geometry's case of 33.2~kJ released over 30~s with a 500~mW\,K$^{-1}$ busbar under the 200 kinetic draws and the 30 thermal draws and recorded when each node's self-heating first reached 0.15~K\,min$^{-1}$: over the kinetic draws the hot spot's onset came 6.01 to 17.13~s after the trigger's start, at 6.25, 10.23, 16.04 and 17.02~s under the four that step 11 took and at 11.48~s under the posterior mean, with a standard deviation of 2.83~s over the kinetic draws and 0.076~s over the thermal draws. A failure counts as foreseen when the predictions file recorded it, before the runs, as expected, as fixed by earlier records or an energy balance, or as following from a condition it named that then held; a hedged ``may fail'' does not count. Where the foreseeing is in the file's analysis rather than in its expectation, the registry quotes that analysis beside the expectation. By that rule 7 of the 20 failures that Table~\ref{tab:s-registry} lists were foreseen. In step 13, P7's first clause, the ignition temperature of the posterior-mean kinetics, was known when the predictions were written and is not scored; P7's verdict is that of its second clause.

Before any case was run, the posterior-mean kinetics were put through the accelerating-rate calorimetry that the kinetic source simulated for its reference cell~\cite{koenig2025uncertainty}: its cell of 720~g and 1100~J\,kg$^{-1}$\,K$^{-1}$, 0.04~m$^2$ and 25~W\,m$^{-2}$\,K$^{-1}$, heated from 25\degC{} in steps of 5~K, each followed by a wait of 30~min and a seek of 10~min for self-heating of 0.15~K\,min$^{-1}$ (the chamber ramping at 1~K\,min$^{-1}$ and the seek's length fixed here, which the source does not give). The exotherm onset is the cell's temperature when a seek first finds that rate, 179.3\degC, inside the source's band of two standard deviations about its main cluster's mean (183.5\degC), 169.0 to 198.0\degC, as read from its figure; the self-heating reaches 1~K\,min$^{-1}$, the source's internal short, at 199.8\degC, below the band of that event, 255.7 to 302.1\degC, read the same way. The 200 kinetic draws went through the same emulation before any case ran. A draw whose self-heating reaches 1~K\,min$^{-1}$ within 4,000~min of the test's start is in the main cluster, the rule the source's own appendix suggests for its figure; 164 of the 200 are, against the source's 76\,\%. The check recorded for them before it ran, that the central 95\,\% of the exotherm onsets of the draws in the source's main cluster lie within 5~K of the band's edges and their mean within 5~K of 183.5\degC, failed: the 164 such draws gave 146.3 to 191.4\degC, with a mean of 178.8\degC. After the failure the check was amended before any case ran, so that it is not a blind check: it took the band's own statistic, the draws' mean and two standard deviations, 161.9 to 195.7\degC, judged about the emulation's onset under the posterior mean, 179.3\degC, as the band is about its mean. So amended it passed; the registered failure stands beside it.

Step 5's 5,248 runs had the cans in contact and the trigger releasing at most 33.2~kJ over 10 or 30~s; at its central case, the central contact and tab with the posterior-mean kinetics, the hottest point peaked at 191.5 to 192.3\degC{} and the lumped neighbour at 159.4 to 162.1\degC{} over the two releases and connections. Step 5's P5 gives the measured mean times of the two connections as 8~s apart; by Table~1 of the main text they are 8.7~s apart. They were: 36 runs at the posterior-mean kinetics over every combination of the three contacts, the three tab conductances (5, 6.8 and 10~mW\,K$^{-1}$), the two releases and the two connections; 8 at the central contact and tab and the 30~s release, with the ambient at 20 and 30\degC{} and with the body energy at the smallest and the median of the Databank's ten nail tests of this cell, 18.67 and 25.88~kJ; 4 with every mesh spacing, or every step limit, halved; 400, the 200 kinetic draws at the central case and the 30~s release in both connections; and the factorial of 4,800, the 200 draws by the three contacts, the tab bracket (5 and 10~mW\,K$^{-1}$), the two releases and the two connections.

\section{The registry of predictions}
\label{sec:s-registry}

The study registered 109 predictions in all, of which 65 held and 44 failed; Table~\ref{tab:s-registry} lists the 63 that the runs reported in the main text and here bear on, 43 of which held and 20 failed, and every prediction, with its predictions file and its verdict, is in the public archive at doi:10.5281/zenodo.23279455~\cite{satte2026repo}. The other 46 are those of steps 1 to 4 and 6 and those of steps 7 to 11 and 13 whose runs neither document reports, such as step 7's runs with an aged SEI or step 8's that heated the parallel neighbour from inside by a share of its stored energy, standing for its discharge through the trigger; Table~\ref{tab:s-steps} gives each step's count, with how many of its listed and of its other predictions held and failed, and a listed prediction's expectation is given without the reasons in it that rest on such runs. Each row gives the step, the prediction's label, its statement as the predictions file gives it (shortened where long, with the test names, the word step and a few of the records' terms written as here), the expectation recorded before the runs, and the verdict. The expectations quote the files: their numbers are what the files recorded before the runs, from hand estimates (step 5; Section~\ref{sec:s-key}), from the reduced models' values registered beside them (steps 13 and 15) and from earlier steps' results that the tables here give, and are not results of this paper.

Step 9's P1, and what step 12's predictions file recorded for its P1 and P2, rest on runs that no other table gives; those runs are given here and in Table~\ref{tab:s-quoted}. Step 9's P1 rests on 83 of that step's runs with the 21.4~s pulse, those at or above the lower end of each propagated test's local threshold: in each test the latest crossing of 255.7\degC{} by the hottest point came 13.8 to 21.1~s after the trigger's start, against measured times of 147 to 342~s. Step 12's predictions file recorded P1 as fixed and P2 as close to fixed by six runs of step 6 with the trigger releasing 33.2~kJ over 21.39~s under the posterior mean: at the central contact and tabs, P1's two configurations, the hottest point peaked at 192.16\degC{} in series and 191.90\degC{} in parallel, as it did again in step 12 (Section~3.1 of the main text), and at the high contact with the low and the high tabs at 226.14 to 226.75\degC. Step 11's P4 was recorded as failing before any of its runs on step 7's runs at the low contact, whose far-face peaks Table~\ref{tab:s-farface} gives, and step 12's P3 and P4 as fixed by step 9's needs, which Tables~\ref{tab:s-thresholds} and~\ref{tab:s-needs} give. Step 15's P1 holds partly by construction: its searches opened 5\,\% either side of the reduced model's value, so that a bracket closed by its first two runs, as three of the four were, has that value as its midpoint; its file did not record it as fixed, so it is not among the seven of Section~2.7 of the main text. Step 16's P7 holds by construction in the same way, as its file recorded, PT-Exp1's bracket having closed at its first two runs, and its file recorded P4 as fixed, or nearly, by earlier records in seven of the eight tests and P11 in most of its cases. These three and step 12's P1 to P4 are the seven held predictions whose expectations in Table~\ref{tab:s-registry} record them as fixed, or nearly.

\IfFileExists{tables/s2_registry.tex}{{\sloppy\scriptsize\setlength{\tabcolsep}{3pt}
\begin{longtable}{>{\raggedright\arraybackslash}p{0.05\textwidth}>{\raggedright\arraybackslash}p{0.065\textwidth}>{\raggedright\arraybackslash}p{0.425\textwidth}>{\raggedright\arraybackslash}p{0.24\textwidth}>{\raggedright\arraybackslash}p{0.1\textwidth}}
\caption{The 63 predictions whose runs the main text or this supplement reports, of the 109 the study registered: 43 held and 20 failed, 7 of the failures foreseen. Each is given as its predictions file states it (shortened where long), with the expectation recorded before the runs and the verdict. Labels are defined in Section~\ref{sec:s-key}.}\label{tab:s-registry}\\
\toprule
Step & & Prediction & Expectation & Verdict \\
\midrule
\endfirsthead
\multicolumn{5}{l}{\textit{Continued}} \\
\toprule
Step & & Prediction & Expectation & Verdict \\
\midrule
\endhead
5 & P1 & At the central contact conductance with the posterior-mean kinetics, the resolved neighbour propagates in series and parallel, its time to propagation within the measured 147 to 342 s. & Study's expectation: May fail: by the estimates above the hot spot stays below 255.7\degC, or, if the line of contact rises further, reaches it within the first minute. & failed \\
5 & P2 & In both connections, series and parallel, the lumped neighbour propagates more than 20 s later than the resolved one, or not at all within 900 s. & Study's expectation: Held, most likely as "not at all": its passive peak should lie near 145 to 170\degC. & held \\
5 & P3 & At the resolved model's propagation time, its mid-surface temperature lies between 150 and 255.7\,\textdegree C: the hot spot leads the surface a thermocouple sees. & Study's expectation: Toward the trigger, held if propagation occurs (that face is the line of contact); away from it, may fail below 150\degC. & failed \\
5 & P4 & The resolved model's mean mid-surface heating rate over the 60 s before propagation lies between 0.2 and 0.8\,\textdegree C/s (measured about 0.4). & Study's expectation: Uncertain. & failed \\
5 & P5 & The predicted series and parallel times to propagation differ by less than 30 s (the measured means differ by 8 s). & Study's expectation: Held if both propagate: the tabs carry 4 to 8 \% of the contact path. & failed \\
5 & P6 & The share of the 4,800 factorial runs in which the resolved neighbour propagates lies between 30 and 90 \% (measured 9 of 16, 56 \%). & Study's expectation: May fail low: the model lacks the vent, ejecta and flame heat, and the estimates above fall short of 255.7\degC{} at the central contact. & failed \\
5 & P7 & The contact and tab conductance range carries more of the spread in predicted time to propagation than the kinetic posterior (standard deviation ratio above 1), reversing step 3's onset finding. & Study's expectation: Held if enough runs propagate to give both spreads: the contact range changes the path 4.6-fold. & failed \\
7 & P1 & At the central contact with per-test triggers, C1 (hot spot at 255.7\degC) predicts propagation in fewer than 9 of the 16 tests. & Plan's expectation: Likely; step 6 ignited the hot spot only at the Databank's total energy and, at the median peak, only at the high contact. Study's expectation: Holds. & held \\
8 & P1 & Of the nine propagated tests, $f^*$(C1) is 0 in the six step 7 hits and below 0.3 in the other three. & Plan's expectation: Likely. Study's expectation: Holds. & held \\
8 & P2 & Of the seven non-propagated tests, at least three, including the two step 7 false alarms at 0, have $f^*$(C1) below 0.3, so the propagated and non-propagated tests' $f^*$(C1) ranges overlap. & Plan's expectation: Likely. Study's expectation: Holds. & held \\
8 & P3 & $f^*$(C2) lies between 0.05 and 0.5 in at least twelve of the sixteen tests. & Plan's expectation: Uncertain. Study's expectation: Fails. & failed, foreseen \\
8 & P6 & Within each connection, $f^*$(C1) and $f^*$(C2) never increase with trigger energy (no inversion). & Plan's expectation: Likely; a check on the mapping. Study's expectation: Holds, with ties. & held \\
9 & P1 & With the 21.39 s delivery, C1 reproduces no propagated test's measured time within 30 s at any fraction at or above $f^*$'s lower end. & Plan's expectation: Likely. Study's expectation: Holds. & held \\
9 & P2 & With the 21.39 s delivery, C2 reproduces at least six of nine propagation times within 30 s, at a fraction between $f^*$'s lower end and 0.1 above its upper end. & Plan's expectation: Uncertain. Study's expectation: Holds. & held \\
9 & P3 & Every C2 time-matching fraction found in P2 lies no more than 0.02 above the upper end of $f^*$(C2). & Plan's expectation: Uncertain. Study's expectation: Uncertain; it turns on PP-Exp1. & held \\
9 & P4 & For the five low-maximum non-propagated tests, the peak-matching fraction at the thermocouple node lies below both $f^*$(C1) and $f^*$(C2). & Plan's expectation: Likely. Study's expectation: Fails, on ST-Exp1 and ST-Exp2, whatever the runs show. & failed, foreseen \\
9 & P5 & The peak-matching fractions (five non-propagated tests) and the C2 time-matching fractions (nine propagated tests) separate: the largest of the former is below the smallest of the latter. & Plan's expectation: Uncertain; this is the step's question. Study's expectation: Fails. & failed, foreseen \\
9 & P7 & The bracket on the ejected-to-body energy ratio rescales every time- and peak-matching fraction by 0.75 to 1.40 and changes P5's verdict in neither direction. & Plan's expectation: Likely. Study's expectation: Holds. & held \\
10 & P1 & Before propagation, the model under the 21.39 s delivery, at one end of its threshold bracket, fits PT-Exp4's measured history at the step 9 (far-face) node within 25 K root-mean-square. & Plan's expectation: uncertain. & held \\
10 & P6 & Held out, no predictor beats one probability for every test by more than 0.1 in mean log score in either split direction, and no permutation p-value is below 0.05. & Plan's expectation: likely. & held \\
10 & P8 & The survey finds at least one dataset with per-test outcomes and times for cells in direct contact, and none whose cell this study's parameters cover without an invented value. & Plan's expectation: uncertain. & failed \\
11 & P1 & For each of the four tests, $f^*$(C2) under every kinetic draw lies within a factor of two of its posterior-mean value, or within 0.025 of it where that is 0. & Plan's expectation: likely. & failed \\
11 & P2 & For each of the four tests, $f^*$(C1) under every kinetic draw lies within a factor of two of its posterior-mean value, or within 0.025 of it where that is 0. & Plan's expectation: uncertain; the hot spot's onset is local and the kinetics set it. & held \\
11 & P3 & Under each of the eight draws, the range of the C2 threshold energy across the four tests exceeds the largest within-test band of that energy across the draws. & Plan's expectation: likely. & held \\
11 & P4 & In the four tests with zero peak-matching fraction, the far face with no delivery exceeds the measured maximum under at least three of four thermal draws and at both contacts. & Plan's expectation: uncertain; the low contact and the smallest conductivity across the layers are the two that could remove it; the predictions file's analysis: “P4, which asks that it do so at both contacts, fails before any step 11 run” & failed, foreseen \\
11 & P5 & In each of the four tests the band of $f^*$(C2) across the four thermal draws is narrower than its band across the four kinetic draws. & Plan's expectation: uncertain. & failed \\
12 & P1 & With $E_\mathrm{b}$ capped at 33.2 kJ, conduction and radiation alone meet C1 in none of the eight tests at the central contact. & Plan's expectation: likely (step 5 found no ignition at 33.2 kJ with the two-cell study's trigger); the predictions file records it as fixed by step 6's records (Section~\ref{sec:s-registry}). & held \\
12 & P2 & At the high contact the same holds. & Plan's expectation: uncertain; the predictions file records it as close to fixed (step 6, with other tabs; Section~\ref{sec:s-registry}). & held \\
12 & P3 & With $E_\mathrm{e}$ fixed at the Databank's median, the delivered energies $E^*$ across the nine propagated tests still span at least tenfold. & Plan's expectation: likely; the predictions file records it as fixed by step 9's records. & held \\
12 & P4 & With $E_\mathrm{e}$ fixed, the propagated and non-propagated tests' fractions still overlap (the largest $f_\mathrm{max}$ above the smallest $f_\mathrm{t}$(C2)). & Plan's expectation: likely; the predictions file records it as fixed by step 9's records. & held \\
13 & P1 & For PT-Exp1 with a 5 s triangular delivery over the facing half, the resolved $f^*$(C1) lies within 10 \% of the skin model's closed form (0.07701). & likely & held \\
13 & P2 & At 40 s (PT-Exp1, facing half), the cross-section conduction model (0.21775) lies within 5 \% of the resolved bracket's midpoint. & uncertain & held \\
13 & P3 & At 120 s (PT-Exp1, facing half), the resolved threshold lies more than 5 \% below the cross-section conduction model's (0.41450): local ignition. & uncertain & failed \\
13 & P4 & For ST-Exp3 and PT-Exp3, the cross-section model with the trigger's conduction lies within 10 \% of the resolved midpoints at 10 and 60 s, and gives the direction of change from 21.39 s in all four cases. & uncertain & held \\
13 & P7 & The ignition temperature of the posterior-mean kinetics lies between 215 and 240\degC{} (known when registered, so not scored), and the central 95 \% of the 200 draws spans less than 40 K (scored). & likely for the mean, uncertain for the span & held \\
13 & P8 & The two-body model's whole-cell threshold lies within the larger of 10 \% and 300 J of the lumped model's midpoint in at least 28 of the 32 low and high contact searches. & uncertain & held \\
13 & P9 & For PT-Exp1, the ratio $f^*$(C2)/$f^*$(C1) falls from 21.39 to 120 to 360 s and is below 1.5 at 360 s. & likely for the fall, uncertain for the level & held \\
13 & P10 & For PP-Exp3, ST-Exp2 and PT-Exp3, the lumped model's needed delivery falls by 150 to 300 J per kJ of the trigger's body energy across the 4 kJ between shifted runs. & likely for the first two, uncertain for PT-Exp3 & failed \\
14 & P1 & The upper bound of radiation from a cap-end vent to the facing half is below 2 \% of 76.74 kJ. & the plan's line for it says likely; the predictions file's analysis, written before any bound: with no source placing the flame, the bound takes the largest share of an emitter's directions that strike the facing half, at least a quarter at the half's end edge, so P1 is expected to fail & failed, foreseen \\
14 & P2 & The upper bounds of the cap-end paths that are determinable sum to less than 5.96 kJ, the largest need among the seven matched times. & uncertain; the predictions file's analysis: with no source for the share of the vent gas that passes the neighbour, the whole of it counts, and its flux over 21.39 s is expected to exceed 5.96 kJ by a wide margin, so P2 is expected to fail & failed, foreseen \\
14 & P3 & They sum to more than 0.26 kJ, the smallest need. & likely & held \\
14 & P4 & The upper bound of a breach's jet facing the neighbour exceeds 5.96 kJ delivered within 21.39 s. & likely & held \\
14 & P5 & At least one of the six paths is not determinable from published inputs for this cell. & likely & held \\
14 & P6 & For the parallel tests the circuit's upper bound on the neighbour's own heating exceeds 1 kJ. & uncertain; the predictions file's analysis: the bound is the neighbour's stored energy, 65,520 J, so P6 is expected to hold & held \\
15 & P1 & The reactive cross-section's $E^*$(C8) lies within 15 \% of the midpoint of the three-dimensional bracket in at least three of four tests. & uncertain; the reduced model's values are the searches' centres & held \\
15 & P2 & For PT-Exp4 with the 21.39 s delivery, $E^*$(C8) lies between 6 and 10 kJ. & likely; the reduced model's value, registered with it: 7.06 kJ & held \\
15 & P3 & In at least three of the four tests $E^*$(C8) is at least 1.25 times the whole-cell clock's threshold $E^*$(C2). & likely for the three strong triggers, uncertain for ST-Exp3; the reduced model's values: 1.49, 1.72, above zero in SP-Exp3, 0.85 & held \\
15 & P4 & For PT-Exp4, in the run at the upper end of the bracket, the mean crosses 255.7\degC{} between 150 and 400 s after the trigger's start. & likely; the reduced model's value: 310.6 s at the upper end of its own bracket & held \\
15 & P5 & In that run the far face reads more than 15 K above 197.3\degC, the reading from which the measured neighbour jumped, when its own jump begins. & likely; the reduced model's value: 240.1\,\textdegree C & held \\
15 & P6 & The constant power after the 5.0 kJ pulse that brings the cell's own runaway to 271.5 s lies between 5 and 25 W. & likely; the reduced model's value: 12.5 W & held \\
15 & P7 & Under none of four kinetic draws and not with the aged SEI does the resolved cell run away at PT-Exp4's 5,045 J, and the least shift of the rate laws that makes it run away is 10 K or more. & likely; the reduced model's values: no runaway in any of the five runs, the least shift 20.0 K & held \\
15 & P8 & PT-Exp4's run with 6,209 J over 240 s, carried to 900 s, runs away. & likely; the reduced model's value, registered with it: no runaway, the mean peaking at 220.5\,\textdegree C & failed \\
16 & P1 & With the 21.4 s pulse, the reactive cross-section's need lies within 10 \% of the three-dimensional need in each of the eight strong-trigger tests that propagated. & holds, the cross-section's need 2 to 9 \% above, since its search stopped at the first run within 10 s of the measured time and the three-dimensional runaway came earlier near the thresholds & held \\
16 & P2 & With the pulse, the cross-section's need lies above the three-dimensional need in each of the eight tests. & holds, for the reasons of P1; least sure in PP-Exp4 and ST-Exp4, whose needs lie near their thresholds & failed \\
16 & P3 & With the pulse, every three-dimensional need lies within 4.1 to 8.1 kJ. & holds; the cross-section's runs put the needs at 4.31 to 7.46 kJ & held \\
16 & P4 & With the pulse, the resolved need exceeds the lumped need in each of the eight tests. & holds; fixed, or nearly, by earlier records in seven of the eight, where a three-dimensional run at or above the lumped need already ran away late or not at all & held \\
16 & P5 & With the pulse, the resolved need exceeds the lumped need by 1.0 to 5.0 kJ in each of the eight tests. & holds, by 1.4 to 4.1 kJ on the cross-section's runs & held \\
16 & P6 & After PT-Exp1's weak trigger the resolved threshold with the pulse lies below the lumped threshold, in three dimensions. & holds; the cross-section's threshold is 0.84 times the lumped one, and the three-dimensional thresholds lay within 15 \% of the cross-section's in the four tests searched & held \\
16 & P7 & PT-Exp1's ratio of the resolved to the lumped threshold, by the brackets' midpoints, lies between 0.75 and 0.95. & holds, at 0.84, fixed by construction if the search's first two runs close its bracket & held \\
16 & P8 & Refining the mesh by 1.5 and halving every step limit moves a need by at most one bracket width. & holds, most likely with both ends of the bracket keeping their outcome & held \\
16 & P9 & That refinement moves the runaway time by less than 3 \% at both ends of the need's bracket. & holds, at about +1.5 \%, as the earlier refinement at PT-Exp4's threshold & held \\
16 & P10 & With constant power, the cross-section's needs lie within 10 \% of the three-dimensional needs found. & fails in at least one search, since with constant power the cross-section's runaway came up to 34.5 s after the three-dimensional model's & failed, foreseen \\
16 & P11 & With constant power, the resolved need exceeds the lumped need in every three-dimensional search closed. & holds; fixed, or nearly, by earlier records in 21 of the 24 cases & held \\
\bottomrule
\end{longtable}
}
}{}

\section{The thresholds and needs of every test}
\label{sec:s-thresholds}

Table~\ref{tab:s-thresholds} gives, test by test, the brackets that Fig.~3 of the main text draws and the needs with the pulse that Fig.~2a draws; Table~\ref{tab:s-needs} gives the needs under every shape. The local and lumped brackets of most tests come from searches over fractions of each test's ejected energy that halved the bracket to 0.0125 of it; the brackets of PT-Exp1, PP-Exp3, ST-Exp2 and SP-Exp3 were narrowed further with one ejected energy for every test, the local brackets of PT-Exp1, ST-Exp3 and PT-Exp3 and the lumped bracket of PT-Exp1 again in step 13 (Table~\ref{tab:s-reduced}), and the lumped brackets of PT-Exp4 and ST-Exp3 by step 12's search for the measured time, which ended at the threshold. PT-Exp4's had first been set by step 9's search at 4.66 to 5.04~kJ (3.75 to 4.06\,\% of its ejected energy), and Sections~3.4 and 3.5 of the main text use its upper end. The cross-section's brackets are at most 1\,\% of their midpoints; the three-dimensional threshold searches started 5\,\% either side of the cross-section's value, and in PT-Exp4 and ST-Exp4 the two runs that bracket the need with the pulse bracket the threshold too, since the lower does not run away within 900~s (Section~\ref{sec:s-step16}).

\IfFileExists{tables/s3_thresholds.tex}{{\footnotesize\setlength{\tabcolsep}{3pt}
\begin{longtable}{p{0.1\textwidth}p{0.13\textwidth}p{0.15\textwidth}p{0.15\textwidth}p{0.15\textwidth}p{0.18\textwidth}}
\caption{The thresholds of the sixteen tests with the 21.4~s pulse at the central contact (kJ; a bracket is the largest energy run that does not meet the criterion and the smallest that does), and, for the nine propagated tests, their needs: the least energy with which the lumped neighbour and the resolved neighbour (by the reactive cross-section, and in the eight tests with strong triggers bracketed in three dimensions) run away within 10~s of the measured time, counted as delivered by the run's own runaway. An energy marked $^{t}$ is the least energy run whose runaway comes before the measured time when the next lower run, 0.5\,\% of the test's ejected energy below, gives none within 900~s: no run reproduces the measured time, and the value, the higher, is the threshold to within that resolution and can exceed the bracket beside it. $^{n}$: a three-dimensional threshold bracketed by the two runs that bracket the need, whose lower end does not run away within 900~s (Section~\ref{sec:s-step16}); PT-Exp4's threshold search gave (6.71, 7.42]. Supplementary Table~\ref{tab:s-needs} gives the other delivery shapes.}\label{tab:s-thresholds}\\
\toprule
Test & Hottest point & Lumped & Resolved, cross-section & Resolved, three-dimensional & Needs: lumped; cross-section; three-dimensional \\
\midrule
\endfirsthead
\multicolumn{6}{l}{\textit{Continued}} \\
\toprule
Test & Hottest point & Lumped & Resolved, cross-section & Resolved, three-dimensional & Needs: lumped; cross-section; three-dimensional \\
\midrule
\endhead
PT-Exp1 & (4.36, 4.53] & (12.85, 12.89] & (10.75, 10.81] & (10.24, 11.32] &  \\
PT-Exp2 & (0.00, 1.68] & (1.68, 3.37] & (6.33, 6.38] &  &  \\
PT-Exp3 & (2.61, 2.68] & (7.36, 8.58] & (8.56, 8.62] &  &  \\
PT-Exp4 & (0.00, 1.55] & (4.66, 4.85] & (7.04, 7.08] & (7.06, 7.15]$^{n}$ & 4.85$^{t}$; 7.44$^{t}$; (7.06, 7.15] \\
PP-Exp1 & 0 & 0 & (4.53, 4.56] &  & 2.97; 6.04; (5.46, 5.65] \\
PP-Exp2 & (0.00, 1.59] & (3.18, 4.77] & (6.83, 6.88] & (6.51, 7.20] & 5.96; 7.65; (7.42, 7.54] \\
PP-Exp3 & (0.00, 0.85] & (2.56, 3.41] & (6.21, 6.25] &  &  \\
PP-Exp4 & 0 & (0.00, 1.82] & (5.58, 5.62] &  & 1.82; 6.05$^{t}$; (5.52, 5.65] \\
ST-Exp1 & 0 & (0.00, 1.98] & (4.42, 4.44] &  &  \\
ST-Exp2 & 0 & (1.80, 2.70] & (5.58, 5.62] &  &  \\
ST-Exp3 & (4.04, 4.20] & (11.95, 12.11] & (10.25, 10.31] & (9.77, 10.80] & 12.11$^{t}$; 10.31$^{t}$ \\
ST-Exp4 & 0 & (0.00, 1.89] & (5.03, 5.06] & (5.32, 5.51]$^{n}$ & 1.18; 5.51$^{t}$; (5.32, 5.51] \\
SP-Exp1 & 0 & (1.82, 3.63] & (5.50, 5.54] &  & 2.73; 5.96; (5.96, 6.15] \\
SP-Exp2 & 0 & (0.00, 1.89] & (5.06, 5.08] &  & 3.78; 6.50; (6.33, 6.50] \\
SP-Exp3 & 0 & 0 & (3.52, 3.54] & (3.70, 4.06] & 0.26; 4.58; (4.45, 4.58] \\
SP-Exp4 & (0.00, 1.54] & (4.61, 6.15] & (7.08, 7.12] &  &  \\
\bottomrule
\end{longtable}
}
}{}

\section{The checks}
\label{sec:s-checks}

Every run's energy account closes: the heat taken in through the boundary paths, the reaction heat and the delivered energy equal the change in stored heat to the residuals of Table~\ref{tab:s-checks}, against an acceptance of $10^{-9}$ of the heat taken in, and every delivery equals its pulse's integral up to the run's end. Each step that changed the configuration of the resolved neighbour halved every mesh spacing, and separately every step limit, on a run near its threshold; step 15 refined the mesh by a factor of 1.5 instead, at both ends of PT-Exp4's threshold search's bracket, where neither change altered the outcome; and step 16 refined the mesh by 1.5 and halved every step limit together at both ends of PP-Exp4's need with the pulse, which delayed the runaway by 1.1 and 2.7\,\% and left both ends in their class. Every registered number was recomputed from the run records by code separate from the code that produced it, and before each step's results were written up each number was recomputed once more from the run records and each verdict read again against them under its registered rule.

\begin{table}[!htbp]
\centering
\footnotesize
\caption{The checks of the steps whose runs the main text or this supplement reports, in full or in part: the largest energy residual of each step's runs, as a fraction of the heat taken in, and the change in the hottest point's peak rise (or a runaway time) with the mesh spacing halved (in steps 15 and 16, refined by a factor of 1.5) and with every step limit halved (in step 16, together with the mesh).}
\label{tab:s-checks}
\setlength{\tabcolsep}{4pt}
\begin{tabular}{rrlll}
\toprule
Step & Runs & Largest energy residual & Mesh refined & Step limits halved \\
\midrule
5 & 5,248 & $8.5\times10^{-12}$ & 0.27 and 0.31\,\% & 0.0035 and 0.0040\,\% \\
6 & 506 & $8.5\times10^{-12}$ & 0.28 and 0.32\,\% & 0.0033 and 0.0037\,\% \\
7 & 186 & $4.34\times10^{-12}$ & 0.15\,\% & 0.0008\,\% \\
8 & 230 & $1.97\times10^{-11}$ & 0.52\,\% & 0.0021\,\% \\
9 & 670 & $6.8\times10^{-12}$ & 0.32\,\% & 0.0006\,\% \\
10 & 30 & $3.3\times10^{-12}$ & \multicolumn{2}{l}{the configuration of step 9, unchanged} \\
11 & 216 & $6.1\times10^{-12}$ & \multicolumn{2}{l}{the configuration of step 9, unchanged} \\
12 & 10 & $1.2\times10^{-12}$ & \multicolumn{2}{l}{the configuration of step 9, unchanged} \\
13 & 389 & $2.99\times10^{-12}$ & \multicolumn{2}{l}{the configuration of step 9, unchanged} \\
15 & 27 & $1.97\times10^{-12}$ & 1.5\,\% (runaway time) & 0.01\,\% (runaway time) \\
16 & 116 & $2.0\times10^{-12}$ & \multicolumn{2}{l}{1.1 and 2.7\,\% (runaway time, both)} \\
\bottomrule
\end{tabular}
\end{table}

\section{The inputs of the bounds on what each path can deliver}
\label{sec:s-inputs}

Each input of the bounds of Table~4 of the main text was read at its source, with the page used, and labelled as measured on this cell, in this rig, on another cell, or as a relation; a value measured on another cell widens a range and sets no value for this one (Table~\ref{tab:s-inputs}). The sources come from a survey whose searches were written and committed before any was run, a search added later being recorded with its reason; the survey, its searches and what each found are in the archive. It looked, for cells of this format and chemistry and otherwise for the nearest, for the vent gas's amount, composition and temperature, the split of the ejecta and the solids' temperature, the jet, the gas's heat of burning, radiant fractions, any heat flux or deposited energy measured where a neighbour stands, the relations for a jet striking a cylinder and for flow past one, the rig's holder and enclosure, the resistances of the parallel tests, where a sidewall breaches and the paint's emissivity. It found no heat flux or deposited energy measured at a neighbouring cylindrical cell's place and no radiant fraction for a battery's vent flame. Where no source gives a fraction, it enters at 1, the largest it can be, so each upper bound is the most that the published inputs surveyed allow under the settings of Table~4 of the main text; the gas's rests on a correlation for convection, and the flame temperature is a module's, cited at second hand, so these bounds are estimates of that most rather than strict limits. Radiation was bounded by tracing two million directions from each of 181 emitter points beyond the trigger's end: six on the vent's axis, from the end face out to ten radii; 147 on a polar grid over the end face, at radii of a quarter to a whole radius in steps of a quarter, its centre among them, and every 30\textdegree, at 0.01, 0.1 and 1 radius beyond its plane; and 28 approaching the facing half's end edge, every 15\textdegree{} from 0 to 90\textdegree{} from the line of contact and 0.01 to 0.3 radius outside the can. On the vent's axis the largest share of an emitter's directions that strike the facing half is about 0.0038, from a finer scan in steps of 0.05 radius from 0.5 to 1.5 radii with four million directions per point, and in front of the trigger's end face, near the line of contact, the share rises toward a quarter. With the holder able to reflect, the geometry stops mattering, and radiation is bounded by a black body at the flame temperature filling the facing half's view. A narrow jet along the can's normal meets the neighbour only if it leaves within 30\textdegree{} of the line of centres, where its path passes within one radius of the neighbour's axis, a sixth of the circumference; the cells lie horizontally, the trigger above~\cite{fransson2024exploring}, so melt from a breach on its lower half falls onto the facing half.

\begin{table}[!htbp]
\centering
\footnotesize
\caption{The inputs of the bounds. Whose: measured on this cell (the LG M50), in this rig, on another cell, or a relation.}
\label{tab:s-inputs}
\renewcommand{\arraystretch}{0.9}
\begin{tabular}{p{0.31\textwidth}p{0.25\textwidth}p{0.2\textwidth}p{0.11\textwidth}}
\toprule
Input & Value & Source & Whose \\
\midrule
Largest energy in the ejecta, ten nail tests & 89.34~kJ (median 76.74~kJ) & \cite{finegan2024databank} & this cell \\
Largest energy released in all, ten nail tests & 110.8~kJ & \cite{finegan2024databank} & this cell \\
Largest heat of burning in open air & 136.5~kJ & \cite{vendra2022numerical}, Table 2 & another cell \\
Largest flame temperature & 1500\degC, a battery module & cited in \cite{chen2020experimental}, p.~2 & another cell \\
Largest vent gas velocity & 398.5~m\,s$^{-1}$ & \cite{gillich2024mechanical} & another cell \\
Shortest runaway or venting & 1.131~s & \cite{gillich2024mechanical} & another cell \\
Longest combustion of a 21700 cell & about 300~s & \cite{chen2023comprehensive}, p.~11 & another cell \\
Sidewall ruptures of the trigger & 3 of the 16 tests & \cite{fransson2024exploring}, Tables 3 and 4 & this cell \\
Sidewall breaches in the Databank's nail tests & none of 10 & \cite{finegan2024databank} & this cell \\
A rupture's jet fire on neighbours & neighbours ran away 2 to 5~s later & \cite{read2024sidewall}, p.~5 & this cell \\
Holder & aluminium frame with 3D-printed inserts, about 2~mm of contact; material not given & \cite{pfaff2023insitu}, p.~194 & this rig \\
Discharge of the neighbour through the trigger & stated for the parallel tests & \cite{fransson2024exploring} & this rig \\
Stored energy at 100\,\% state of charge & 65.52~kJ (18.2~Wh) & \cite{finegan2024databank} & this cell \\
Open-circuit voltage before the test, ten nail tests & 4.164 to 4.189~V & \cite{finegan2024databank} & this cell \\
Direct-current resistance, new cell at 25\degC & $30\pm6$~m$\Omega$ & \cite{lg2017m50}, item 4.2.2 & this cell \\
Resistance's rise as the cell cools & about five times from 45 to 10\degC & \cite{popp2020antemortem}, p.~172 & this cell \\
Convection from a cylinder in cross flow & Churchill and Bernstein's relation & \cite{churchill1977correlating,lienhard2024heat} & relation \\
Properties of hydrogen and air at 1 atm & tabulated to 1000 and 1500~K & \cite[Table A.6]{lienhard2024heat} & relation \\
Radiant fraction, absorptivity, share of gas passing, share carried by a breach & no source; each enters at 1 & & limit \\
\bottomrule
\end{tabular}
\end{table}

The shortest venting, 1.131~s, and the longest combustion, about 300~s, set the window in which a path of the ejecta can act, the second also the time over which the bounds limited by a rate, radiation's and the gas's, are taken in total. The rupture's jet fire, after which Read, Jones and Marco's neighbours ran away 2 to 5~s later, shows that path in this cell; no number of it enters a bound. The resistance's rise as the cell cools, about five times from 45 to 10\degC, shows how far it moves with temperature; no source gives it above 45\degC. The neighbour's direct-current resistance, $30\pm6$~m$\Omega$ at 25\degC, is the manufacturer's specification of the cell~\cite{lg2017m50}; it sets only the rate at which the discharge path could deliver, the largest open-circuit voltage of the Databank's ten tests, 4.189~V, across the least resistance, 24~m$\Omega$, 731~W and so 15.6~kJ within 21.4~s if the resistance kept its value at 25\degC, while the bound in Table~4 of the main text is the whole stored energy.

\section{The one published history of a propagating neighbour}
\label{sec:s-history}

No raw temperature series of the two-cell tests is published. The trigger's release in the model, 21.39~s, is the mean of the two triggers' rise times in the same figure: each the rise from 25\degC, taken as the start of the record, to the trigger's maximum in Fransson et al.'s Table~3 (686\degC{} in PP-Exp3, 630\degC{} in PT-Exp4), over the mean heating rate the figure prints for its Section~I, which its caption calls the region of the trigger's runaway (37.70 and 23.97~K\,s$^{-1}$), which gives 17.53 and 25.24~s. The supplement of Fransson et al.~\cite{fransson2024exploring} draws the thermocouple records of two tests in its Figure~10, PT-Exp4, which propagated, and PP-Exp3, which did not. The figure is a vector drawing, so the series were read out of the curves' paths rather than off pixels, twice, by separate code and separate calibrations, and resampled every 0.5~s from 0 to 500~s. The two readings differ by at most 0.0005~K on PT-Exp4's neighbour curve; the figure's own resolution, half the curves' stroke width, is 2.47~K and 0.45~s, and the recorded uncertainty is the larger. A third reading, from the page rendered at 1200~dpi, lies within the recorded uncertainty plus one pixel of the vector reading. PT-Exp4's record: 32.2\degC{} at 0~s; a spike to 102.0\degC{} at 3.1~s as the trigger runs away, back to 79.9\degC{} at 8.4~s (80.0\degC{} at 8.5~s on the 0.5~s samples); 141.0\degC{} at 45~s; a slow rise to 197.3\degC{} at 271.50~s, where the figure marks the onset (Table~3 of the paper gives 272~s); then the jump. The main text takes the opening 8.5~s as the thermocouple's exposure at the trigger's runaway rather than the can's temperature.

The far face's jump in a run begins at the first time its rate of rise, $[T(t+1\,\mathrm{s})-T(t-1\,\mathrm{s})]/2$~s on a 0.5~s grid, reaches 1~K\,s$^{-1}$ after the rate has first reached 0.5~K\,s$^{-1}$ and then fallen below it, before the run's runaway; where there is no such time, the reading 20~s before the runaway stands in, which no run whose jump the main text quotes needed. The rule was fixed before the runs it reads, as the repository's history shows. Applied to PT-Exp4's record from the end of the figure's Section~I, 24.16~s, it finds the measured jump beginning at 267.5~s and 186.54\degC, 10.72~K below the record's 197.25\degC{} at 271.5~s, which the figure marks and the registered comparison takes as 197.3\degC; read by one rule, the model's jump in PT-Exp4's run at the upper end of its resolved bracket begins 47.6~K above the record's, against that comparison's 36.8~K.

\section{Other datasets}
\label{sec:s-survey}

A survey of published datasets with per-test outcomes and times for cell-to-cell propagation, made by web searches whose queries the archive records and by reading each candidate on its publisher's or repository's page where that page could be fetched, asked whether the model's inputs would cover their cell with no value invented beyond the assumptions the model already makes (Table~2 of the main text; Table~\ref{tab:s-survey}). One other dataset has a cell they cover: Read, Jones and Marco~\cite{read2024sidewall}, seven LG M50 cells 3~mm apart, whose propagated tests followed the trigger's sidewall rupture, a path this model does not represent.

\begin{table}[!htbp]
\centering
\footnotesize
\caption{The datasets the survey found.}
\label{tab:s-survey}
\begin{tabular}{p{0.3\textwidth}p{0.21\textwidth}p{0.17\textwidth}p{0.17\textwidth}}
\toprule
Dataset & Cell & Arrangement & The model's inputs cover the cell \\
\midrule
Fransson et al. 2024~\cite{fransson2024exploring}, M50 tests & LG M50 & two cells in contact & yes \\
Fransson et al. 2024, P42A tests & Molicel P42A & two cells in contact & no \\
Read, Jones and Marco 2024~\cite{read2024sidewall} & LG M50 & seven cells, 3~mm gaps & yes \\
Read et al. 2024~\cite{read2024barrier} & unnamed, M50 specification & 34 cells, 1~mm gaps & not established \\
Allen et al. 2025~\cite{allen2025sleeves} & Molicel P45B & seven cells, 3~mm & no \\
Lamb et al. 2015~\cite{lamb2015failure} & LiCoO$_2$, 18650 and pouch & cells in contact & no \\
Said et al. 2019~\cite{said2019cascading} & LiCoO$_2$, 18650 & arrays & no \\
Hoelle et al. 2023~\cite{hoelle2023simulation} & prismatic & five-cell stack & no \\
Sch\"oberl et al. 2024~\cite{schoberl2024propagation} & prismatic NMC and LFP & braced stack & no \\
Torres-Castro et al. 2020~\cite{torrescastro2020passive} & pouch & stacks & no \\
\bottomrule
\end{tabular}
\end{table}

\section{The needs under every delivery shape}
\label{sec:s-needs}

Table~\ref{tab:s-needs} gives the needs of every test under the four shapes that Fig.~2 of the main text draws: for the nine tests that propagated, the lumped neighbour's and the resolved neighbour's by the reactive cross-section, with the resolved delivery's energy in all and the time of its runaway; for the seven that did not, the energy that puts the far face's peak at the measured maximum. Each search for a need stops when its bracket is narrower than 0.5\,\% of the test's ejected energy $E_\mathrm{e}$. The resolved neighbour's search reads energies above its threshold's upper end by 0, 0.0125, 0.025, 0.05, 0.1 and 0.2 of $E_\mathrm{e}$, in order, up to the first whose runaway comes before the measured time, and then bisects; a dash in the table marks a search whose grid ended without such a run.

\IfFileExists{tables/s7_needs.tex}{{\scriptsize\setlength{\tabcolsep}{3pt}
\begin{longtable}{p{0.08\textwidth}p{0.06\textwidth}p{0.12\textwidth}p{0.09\textwidth}p{0.09\textwidth}p{0.09\textwidth}p{0.08\textwidth}p{0.17\textwidth}}
\caption{The needs of the nine tests that propagated under the four delivery shapes (kJ): the least energy with which the lumped neighbour (Lumped), and the resolved neighbour by the reactive cross-section (Resolved), run away within 10~s of the measured time (Time), counted as delivered by the run's own runaway; for the resolved neighbour, the energy of that delivery in all, over its whole shape, and the time of its runaway in that run. $^{t}$, the least energy run whose runaway comes before the measured time, the next lower run, 0.5\,\% of the test's ejected energy below, giving none within 900~s, so that no run reproduces the measured time and the energy is the threshold to within that resolution; the energy is counted to that run's runaway; $^{b}$, interpolated between the two runs whose runaways bracket the measured time, and counted to the measured time; a dash, no energy on the search's grid, which ran to 0.2 of the test's ejected energy above the threshold's upper end, brought the runaway within 10~s of the measured time (in ST-Exp3 over 360~s it came no earlier than 196.4~s, with 25.6~kJ). The four shapes deliver the same energy by the runaway only where the delivery ends before it. Three-dimensional: the need bracketed in three dimensions (step 16), counted as delivered by 10~s after the measured time; the searches with constant power that the computing allowed are given, and one not finished is marked so. Below, for the seven tests that did not propagate, the energy that puts the far face's peak at the measured maximum, within 5~K, under each shape; $^{a}$, a maximum that may not be the neighbour's surface (exploratory); a dash, a far face that stays below the measured maximum with up to the test's whole ejected energy.}\label{tab:s-needs}\\
\toprule
Test & Time (s) & Delivery & Lumped & Resolved & Resolved, in all & Its runaway (s) & Three-dimensional \\
\midrule
\endfirsthead
\multicolumn{8}{l}{\textit{Continued}} \\
\toprule
Test & Time (s) & Delivery & Lumped & Resolved & Resolved, in all & Its runaway (s) & Three-dimensional \\
\midrule
\endhead
PT-Exp4 & 272 & 21.4~s pulse & 4.85$^{t}$ & 7.44$^{t}$ & 7.44 & 188.5 & (7.06, 7.15] \\
 &  & 120~s power & 3.88$^{t}$ & 7.20$^{t}$ & 7.20 & 200.6 & (6.57, 6.68] \\
 &  & 240~s power & 3.49 & 6.77 & 6.77 & 271.4 & (6.40, 6.50] \\
 &  & 360~s power & 3.76 & 7.28 & 9.39 & 278.9 &  \\
\midrule
PP-Exp1 & 152 & 21.4~s pulse & 2.97 & 6.04 & 6.04 & 156.3 & (5.46, 5.65] \\
 &  & 120~s power & 1.98 & 5.42 & 5.42 & 161.5 & (4.88, 4.99] \\
 &  & 240~s power & 2.42 & 6.00 & 9.49 & 151.7 & (5.12, 5.24] \\
 &  & 360~s power & 2.42 & 5.94 & 13.84 & 154.4 &  \\
\midrule
PP-Exp2 & 147 & 21.4~s pulse & 5.96 & 7.65 & 7.65 & 148.4 & (7.42, 7.54] \\
 &  & 120~s power & 4.77 & 7.74 & 7.74 & 143.7 & (7.09, 7.19] \\
 &  & 240~s power & 4.85 & 8.04 & 12.50 & 154.4 & (7.57, 7.67] \\
 &  & 360~s power & 5.33 & 8.12 & 19.52 & 149.7 &  \\
\midrule
PP-Exp4 & 275 & 21.4~s pulse & 1.82 & 6.05$^{t}$ & 6.05 & 216.9 & (5.52, 5.65] \\
 &  & 120~s power & 1.46$^{b}$ & 5.46$^{t}$ & 5.46 & 228.4 & (4.71, 4.84] \\
 &  & 240~s power & 1.36$^{t}$ & 5.03 & 5.03 & 272.8 & (4.48, 4.66] \\
 &  & 360~s power & 1.39 & 5.40 & 6.91 & 281.3 &  \\
\midrule
ST-Exp3 & 177 & 21.4~s pulse & 12.11$^{t}$ & 10.31$^{t}$ & 10.31 & 114.0 &  \\
 &  & 120~s power & 11.47$^{t}$ & 12.55$^{t}$ & 12.55 & 160.1 &  \\
 &  & 240~s power & 12.26 & 13.84 & 18.77 & 177.0 &  \\
 &  & 360~s power & 12.25 & -- &  &  &  \\
\midrule
ST-Exp4 & 342 & 21.4~s pulse & 1.18 & 5.51$^{t}$ & 5.51 & 231.8 & (5.32, 5.51] \\
 &  & 120~s power & 0.95 & 4.92$^{t}$ & 4.92 & 233.7 & (4.55, 4.67] \\
 &  & 240~s power & 0.95$^{t}$ & 4.48$^{t}$ & 4.48 & 272.6 & (4.10, 4.29] \\
 &  & 360~s power & 1.08$^{t}$ & 4.74 & 4.92 & 347.1 &  \\
\midrule
SP-Exp1 & 223 & 21.4~s pulse & 2.73 & 5.96 & 5.96 & 223.2 & (5.96, 6.15] \\
 &  & 120~s power & 2.17$^{b}$ & 5.40 & 5.40 & 230.1 & (5.30, 5.40] \\
 &  & 240~s power & 2.12 & 5.78 & 6.35 & 218.7 & (5.54, 5.67] \\
 &  & 360~s power & 2.19 & 5.79 & 9.56 & 218.1 &  \\
\midrule
SP-Exp2 & 149 & 21.4~s pulse & 3.78 & 6.50 & 6.50 & 148.6 & (6.33, 6.50] \\
 &  & 120~s power & 2.84 & 5.88 & 5.88 & 158.7 & unfinished \\
 &  & 240~s power & 3.40 & 6.32 & 9.70 & 156.4 & (6.09, 6.18] \\
 &  & 360~s power & 3.04 & 6.50 & 15.84 & 147.7 &  \\
\midrule
SP-Exp3 & 210 & 21.4~s pulse & 0.26 & 4.58 & 4.58 & 209.9 & (4.45, 4.58] \\
 &  & 120~s power & 0.22$^{b}$ & 4.06 & 4.06 & 201.9 & unfinished \\
 &  & 240~s power & 0.23$^{b}$ & 4.18 & 4.80 & 209.2 & (3.88, 4.01] \\
 &  & 360~s power & 0.30 & 4.18 & 7.22 & 208.8 &  \\
\midrule
\multicolumn{8}{l}{\textit{Tests that did not propagate}} \\
Test & Peak (\textdegree C) & Delivery & \multicolumn{5}{l}{Far face's peak at the measured peak (kJ)} \\
\midrule
PT-Exp1 & 71 & 21.4~s pulse & \multicolumn{5}{l}{2.09} \\
 &  & 120~s power & \multicolumn{5}{l}{1.52} \\
 &  & 240~s power & \multicolumn{5}{l}{1.01} \\
 &  & 360~s power & \multicolumn{5}{l}{1.23} \\
PT-Exp2 & 808$^{a}$ & 21.4~s pulse & \multicolumn{5}{l}{75.81} \\
 &  & 120~s power & \multicolumn{5}{l}{96.87} \\
 &  & 240~s power & \multicolumn{5}{l}{--} \\
 &  & 360~s power & \multicolumn{5}{l}{--} \\
PT-Exp3 & 118 & 21.4~s pulse & \multicolumn{5}{l}{0.00} \\
 &  & 120~s power & \multicolumn{5}{l}{0.00} \\
 &  & 240~s power & \multicolumn{5}{l}{0.00} \\
 &  & 360~s power & \multicolumn{5}{l}{0.00} \\
PP-Exp3 & 105 & 21.4~s pulse & \multicolumn{5}{l}{0.00} \\
 &  & 120~s power & \multicolumn{5}{l}{0.00} \\
 &  & 240~s power & \multicolumn{5}{l}{0.00} \\
 &  & 360~s power & \multicolumn{5}{l}{0.00} \\
ST-Exp1 & 128 & 21.4~s pulse & \multicolumn{5}{l}{0.00} \\
 &  & 120~s power & \multicolumn{5}{l}{0.00} \\
 &  & 240~s power & \multicolumn{5}{l}{0.00} \\
 &  & 360~s power & \multicolumn{5}{l}{0.00} \\
ST-Exp2 & 133 & 21.4~s pulse & \multicolumn{5}{l}{0.00} \\
 &  & 120~s power & \multicolumn{5}{l}{0.00} \\
 &  & 240~s power & \multicolumn{5}{l}{0.00} \\
 &  & 360~s power & \multicolumn{5}{l}{0.00} \\
SP-Exp4 & 626$^{a}$ & 21.4~s pulse & \multicolumn{5}{l}{32.67} \\
 &  & 120~s power & \multicolumn{5}{l}{40.35} \\
 &  & 240~s power & \multicolumn{5}{l}{76.86} \\
 &  & 360~s power & \multicolumn{5}{l}{115.30} \\
\bottomrule
\end{longtable}
}
}{}

\section{The local threshold against the reduced models}
\label{sec:s-reduced}

Table~\ref{tab:s-reduced} gives the numbers behind Fig.~4a of the main text and its Section~3.4: the resolved model's local threshold with triangular pulses of 5 to 360~s, against the cross-section conduction model and the skin model, with the hottest point's peak when the reaction is switched off at the bracket's upper end, and PT-Exp1's lumped threshold at three durations. The searches at 5, 40 and 120~s for PT-Exp1 and at 10 and 60~s for ST-Exp3 and PT-Exp3 started at the cross-section model's value with a first step of 3\,\% of it and ended once the bracket was narrower than 4\,\% of its midpoint.

\begin{table}[!htbp]
\centering
\scriptsize
\caption{The local threshold against the reduced models of Section~2.5 of the main text, with triangular pulses over the facing half (kJ): PT-Exp1 by the pulse's duration, and ST-Exp3 and PT-Exp3 at 10, 21.4 and 60~s. Fourier number, $\alpha\tau/a^2$ with the roll's diffusivity across its layers. Skin model: PT-Exp1 only, as Fig.~4a of the main text draws it. Hottest point without the reaction: with the reaction switched off at the upper end of the resolved bracket, the hottest point's peak less 255.7\textdegree C (a negative value falls short). Lumped threshold and the ratio of the lumped to the local threshold's midpoint: PT-Exp1 at the central contact. A blank: no search or run at that duration.}
\label{tab:s-reduced}
\setlength{\tabcolsep}{2.5pt}
\begin{tabular}{p{0.075\textwidth}>{\raggedleft\arraybackslash}p{0.055\textwidth}>{\raggedleft\arraybackslash}p{0.075\textwidth}>{\raggedleft\arraybackslash}p{0.13\textwidth}>{\raggedleft\arraybackslash}p{0.085\textwidth}>{\raggedleft\arraybackslash}p{0.06\textwidth}>{\raggedleft\arraybackslash}p{0.1\textwidth}>{\raggedleft\arraybackslash}p{0.13\textwidth}>{\raggedleft\arraybackslash}p{0.055\textwidth}}
\toprule
Test & Pulse (s) & Fourier number & Local threshold, resolved & Cross-section conduction & Skin model & Hottest point, no reaction (K) & Lumped threshold & Ratio \\
\midrule
PT-Exp1 & 5 & 0.025 & (2.15, 2.21] & 2.15 & 2.15 & 5.5 &  &  \\
 & 10 & 0.050 & (2.79, 3.14] & 2.98 & 2.86 &  &  &  \\
 & 21.4 & 0.107 & (4.36, 4.53] & 4.37 & 3.90 &  & (12.85, 12.89] & 2.90 \\
 & 40 & 0.200 & (5.89, 6.07] & 6.07 & 5.02 & -1.1 &  &  \\
 & 60 & 0.300 & (7.67, 8.01] & 7.71 & 6.05 &  &  &  \\
 & 120 & 0.600 & (11.21, 11.55] & 11.55 & 8.39 & -0.8 & (13.10, 13.30] & 1.16 \\
 & 240 & 1.200 &  & 16.25 & 11.67 &  &  &  \\
 & 360 & 1.800 & (17.06, 17.65] & 19.39 & 14.73 & -18.9 & (16.16, 16.41] & 0.94 \\
\midrule
ST-Exp3 & 10 & 0.050 & (2.97, 3.05] & 2.97 &  & 5.3 &  &  \\
 & 21.4 & 0.107 & (4.04, 4.20] & 4.15 &  &  &  &  \\
 & 60 & 0.300 & (6.43, 6.63] & 6.63 &  & -0.7 &  &  \\
\midrule
PT-Exp3 & 10 & 0.050 & (2.94, 3.03] & 2.94 &  & 5.2 &  &  \\
 & 21.4 & 0.107 & (2.61, 2.68] & 2.59 &  &  &  &  \\
 & 60 & 0.300 & (3.23, 3.33] & 3.23 &  & 1.5 &  &  \\
\bottomrule
\end{tabular}
\end{table}


\section{The held-out scores}
\label{sec:s-heldout}

Table~\ref{tab:s-heldout} gives the held-out test's scores, for each of the three quantities fitted on one connection and scored on the other, in both directions. The need is the lumped neighbour's where a test propagated and the far face's where it did not, so its definition depends on the outcome it is set to predict. In the two tests whose maxima may not be the neighbour's surface the far face reaches the maximum only where the model's neighbour runs away, so their needs, 32.7 and 75.8~kJ, are energies of a runaway; they enter as registered. Each quantity's logistic predictor and the one probability were fitted by likelihood penalised with the Jeffreys prior, which keeps a fit finite where a quantity separates the eight tests and gives the one probability as $(k+0.5)/(n+1)$ for $k$ of $n$ tests propagated; each $p$ is the share of 10,000 permutations of the outcomes whose gain is at least the observed one.

\begin{table}[!htbp]
\centering
\footnotesize
\caption{The held-out test of Section~2.7 of the main text, with the 21.4~s pulse: each quantity fitted as a logistic predictor of the outcome on the eight tests of one connection and scored on the other eight by the mean log score (higher is better), against one probability for every test fitted on the same eight; the gain is the difference, and $p$ the share of 10,000 permutations of the outcomes whose gain is at least the observed one.}
\label{tab:s-heldout}
\setlength{\tabcolsep}{4pt}
\begin{tabular}{p{0.25\textwidth}llrrrr}
\toprule
Quantity & Fitted on & Scored on & One probability & Logistic & Gain & $p$ \\
\midrule
Local threshold & series & parallel & -0.718 & -0.848 & -0.130 & 0.66 \\
 & parallel & series & -0.693 & -0.721 & -0.028 & 0.75 \\
Lumped threshold & series & parallel & -0.718 & -0.802 & -0.084 & 0.59 \\
 & parallel & series & -0.693 & -0.706 & -0.013 & 0.62 \\
Need: lumped where propagated, far face where not & series & parallel & -0.718 & -0.656 & 0.062 & 0.23 \\
 & parallel & series & -0.693 & -0.668 & 0.025 & 0.35 \\
\bottomrule
\end{tabular}
\end{table}


\section{The lumped threshold at the low and the high contact}
\label{sec:s-twobody}

Table~\ref{tab:s-twobody} gives the 32 searches of the lumped threshold at the low and the high contact that the two-body model was registered against; the tolerance is the larger of 10\,\% of the bracket's midpoint and 300~J, as registered with the model; where the thresholds are small the 300~J decides, so that, for one, PT-Exp2's 0.29~kJ at the high contact counts as within a bracket of 0.09 to 0.11~kJ. Nine of the 32 are zero in both. At the low contact the two-body model lies above the bracket's midpoint after every strong trigger, by 0.25 to 0.48~kJ, and below it after the two weak ones, by 0.28 and 0.44~kJ, an offset with the trigger's strength that the tolerance does not show; at the high contact the differences run from $-0.43$ to 0.19~kJ. At the central contact the two-body model was compared, unscored, with the 32 searches made there with triangular pulses of 120 and 360~s, which the lower part of Table~\ref{tab:s-twobody} gives: it lies within the same tolerance in 29, and in the other three, ST-Exp1, ST-Exp4 and SP-Exp2 over 120~s, 7 to 46~J beyond it.

\IfFileExists{tables/s10_twobody.tex}{{\footnotesize\setlength{\tabcolsep}{2.5pt}
\begin{longtable}{llrrrrrl}
\caption{The lumped threshold against the two-body model of Section~2.5 of the main text (kJ): above, with the 21.4~s pulse at the low and the high contact, the 32 searches the model was registered against before them; below, with triangular pulses of 120 and 360~s at the central contact, the 32 searches it was compared with afterwards, unscored. Each row gives the search's bracket, the two-body model's value and its difference from the bracket's midpoint, against the larger of 10\,\% of the midpoint and 300~J, the tolerance registered with the model.}\label{tab:s-twobody}\\
\toprule
Test & Case & $E_\mathrm{b}$ & Lumped bracket & Two-body & Difference & Tolerance & Within \\
\midrule
\endfirsthead
\multicolumn{8}{l}{\textit{Continued}} \\
\toprule
Test & Case & $E_\mathrm{b}$ & Lumped bracket & Two-body & Difference & Tolerance & Within \\
\midrule
\endhead
\multicolumn{8}{l}{\textit{At the low and the high contact, the 21.4~s pulse (registered)}} \\
PT-Exp1 & low & 9.0 & (12.43, 12.62] & 12.25 & -0.28 & 1.25 & yes \\
 & high &  & (12.82, 13.01] & 12.64 & -0.28 & 1.29 & yes \\
PT-Exp2 & low & 43.4 & (8.76, 8.89] & 9.31 & 0.48 & 0.88 & yes \\
 & high &  & (0.09, 0.11] & 0.29 & 0.19 & 0.30 & yes \\
PT-Exp3 & low & 31.6 & (10.74, 10.91] & 11.07 & 0.25 & 1.08 & yes \\
 & high &  & (6.97, 7.07] & 6.97 & -0.05 & 0.70 & yes \\
PT-Exp4 & low & 40.0 & (9.44, 9.58] & 9.88 & 0.37 & 0.95 & yes \\
 & high &  & (2.20, 2.23] & 2.34 & 0.12 & 0.30 & yes \\
PP-Exp1 & low & 51.1 & (7.47, 7.58] & 7.93 & 0.41 & 0.75 & yes \\
 & high &  & 0 & 0.00 & 0.00 & 0.30 & yes \\
PP-Exp2 & low & 41.0 & (9.29, 9.43] & 9.72 & 0.36 & 0.94 & yes \\
 & high &  & (1.60, 1.63] & 1.76 & 0.15 & 0.30 & yes \\
PP-Exp3 & low & 43.9 & (8.68, 8.81] & 9.22 & 0.48 & 0.87 & yes \\
 & high &  & 0 & 0.00 & 0.00 & 0.30 & yes \\
PP-Exp4 & low & 46.9 & (8.18, 8.31] & 8.70 & 0.45 & 0.82 & yes \\
 & high &  & 0 & 0.00 & 0.00 & 0.30 & yes \\
ST-Exp1 & low & 51.1 & (7.83, 7.95] & 8.32 & 0.43 & 0.79 & yes \\
 & high &  & 0 & 0.00 & 0.00 & 0.30 & yes \\
ST-Exp2 & low & 46.5 & (8.57, 8.70] & 9.11 & 0.47 & 0.86 & yes \\
 & high &  & 0 & 0.00 & 0.00 & 0.30 & yes \\
ST-Exp3 & low & 16.6 & (12.13, 12.31] & 11.78 & -0.44 & 1.22 & yes \\
 & high &  & (11.90, 12.08] & 11.56 & -0.43 & 1.20 & yes \\
ST-Exp4 & low & 48.8 & (8.20, 8.33] & 8.71 & 0.45 & 0.83 & yes \\
 & high &  & 0 & 0.00 & 0.00 & 0.30 & yes \\
SP-Exp1 & low & 46.8 & (8.51, 8.64] & 9.05 & 0.47 & 0.86 & yes \\
 & high &  & 0 & 0.00 & 0.00 & 0.30 & yes \\
SP-Exp2 & low & 48.7 & (8.21, 8.34] & 8.73 & 0.45 & 0.83 & yes \\
 & high &  & 0 & 0.00 & 0.00 & 0.30 & yes \\
SP-Exp3 & low & 54.3 & (7.31, 7.43] & 7.78 & 0.41 & 0.74 & yes \\
 & high &  & 0 & 0.00 & 0.00 & 0.30 & yes \\
SP-Exp4 & low & 39.6 & (9.73, 9.88] & 10.19 & 0.38 & 0.98 & yes \\
 & high &  & (2.67, 2.71] & 2.84 & 0.15 & 0.30 & yes \\
\midrule
\multicolumn{8}{l}{\textit{At the central contact, triangular pulses (compared afterwards, unscored)}} \\
PT-Exp1 & 120~s & 9.0 & (13.10, 13.30] & 13.10 & -0.10 & 1.32 & yes \\
 & 360~s &  & (16.16, 16.41] & 16.91 & 0.63 & 1.63 & yes \\
PT-Exp2 & 120~s & 43.4 & (2.53, 2.57] & 2.78 & 0.22 & 0.30 & yes \\
 & 360~s &  & (2.23, 2.27] & 2.45 & 0.20 & 0.30 & yes \\
PT-Exp3 & 120~s & 31.6 & (6.58, 6.68] & 6.58 & -0.05 & 0.66 & yes \\
 & 360~s &  & (6.42, 6.52] & 6.62 & 0.15 & 0.65 & yes \\
PT-Exp4 & 120~s & 40.0 & (3.71, 3.77] & 3.94 & 0.20 & 0.37 & yes \\
 & 360~s &  & (3.37, 3.43] & 3.53 & 0.13 & 0.34 & yes \\
PP-Exp1 & 120~s & 51.1 & 0 & 0.12 & 0.12 & 0.30 & yes \\
 & 360~s &  & 0 & 0.11 & 0.11 & 0.30 & yes \\
PP-Exp2 & 120~s & 41.0 & (3.40, 3.45] & 3.61 & 0.19 & 0.34 & yes \\
 & 360~s &  & (3.08, 3.12] & 3.22 & 0.12 & 0.31 & yes \\
PP-Exp3 & 120~s & 43.9 & (2.34, 2.38] & 2.61 & 0.25 & 0.30 & yes \\
 & 360~s &  & (2.09, 2.13] & 2.30 & 0.19 & 0.30 & yes \\
PP-Exp4 & 120~s & 46.9 & (1.29, 1.31] & 1.58 & 0.28 & 0.30 & yes \\
 & 360~s &  & (1.13, 1.15] & 1.39 & 0.25 & 0.30 & yes \\
ST-Exp1 & 120~s & 51.1 & (0.14, 0.16] & 0.49 & 0.35 & 0.30 & no \\
 & 360~s &  & (0.13, 0.14] & 0.43 & 0.30 & 0.30 & yes \\
ST-Exp2 & 120~s & 46.5 & (1.75, 1.78] & 2.04 & 0.28 & 0.30 & yes \\
 & 360~s &  & (1.53, 1.55] & 1.78 & 0.24 & 0.30 & yes \\
ST-Exp3 & 120~s & 16.6 & (11.04, 11.20] & 10.87 & -0.24 & 1.11 & yes \\
 & 360~s &  & (12.53, 12.72] & 13.12 & 0.49 & 1.26 & yes \\
ST-Exp4 & 120~s & 48.8 & (0.93, 0.95] & 1.26 & 0.32 & 0.30 & no \\
 & 360~s &  & (0.83, 0.84] & 1.10 & 0.26 & 0.30 & yes \\
SP-Exp1 & 120~s & 46.8 & (1.62, 1.65] & 1.93 & 0.29 & 0.30 & yes \\
 & 360~s &  & (1.43, 1.46] & 1.68 & 0.23 & 0.30 & yes \\
SP-Exp2 & 120~s & 48.7 & (0.97, 0.98] & 1.28 & 0.31 & 0.30 & no \\
 & 360~s &  & (0.85, 0.86] & 1.12 & 0.27 & 0.30 & yes \\
SP-Exp3 & 120~s & 54.3 & 0 & 0.00 & 0.00 & 0.30 & yes \\
 & 360~s &  & 0 & 0.00 & 0.00 & 0.30 & yes \\
SP-Exp4 & 120~s & 39.6 & (4.12, 4.19] & 4.32 & 0.16 & 0.42 & yes \\
 & 360~s &  & (3.69, 3.75] & 3.86 & 0.14 & 0.37 & yes \\
\bottomrule
\end{longtable}
}
}{}

\section{The archived runs the reactive cross-section was set against}
\label{sec:s-archived}

Table~\ref{tab:s-archived} lists, by test and delivery, the 151 archived three-dimensional runs at the central contact against which Section~3.7 of the main text sets the reactive cross-section. They are every archived run with the reaction on whose stop at the hottest point's crossing was lifted, whether or not the hottest point then crossed, from earlier steps' searches and their comparison with PT-Exp4's history; the comparison was made after them. The other archived runs stopped once the hottest point reached 255.7\degC, before the mean could. Six repeat another's configuration in a different step, two of them carried only to 332~s, and are counted as they were run; Table~\ref{tab:s-archived-runs} (Section~\ref{sec:s-archived-runs}) lists every run with both models' outcomes and marks the six. Of the 92 in which both models run away, 82 are of the two tests whose maxima may not be the neighbour's surface, PT-Exp2 and SP-Exp4, at deliveries far above any threshold.

\begin{table}[!htbp]
\centering
\footnotesize
\caption{The 151 archived three-dimensional runs at the central contact against which the reactive cross-section was set (Section~3.7 of the main text), by test and delivery: the runs with no delivery, with the 21.4~s pulse and with constant power over 120, 240 and 360~s; the range of energy delivered; the runs in which both models' means reach 255.7\textdegree C; and those held to 332~s, PT-Exp4's measured time plus 60~s, made to compare with its history, every other run going to 900~s. In all 151 the cross-section gave the three-dimensional outcome.}
\label{tab:s-archived}
\setlength{\tabcolsep}{2.5pt}
\begin{tabular}{lrrrrrr@{\hspace{9pt}}lrr}
\toprule
Test & Runs & None & Pulse & 120~s & 240~s & 360~s & Energy (kJ) & Both run away & To 332~s \\
\midrule
PT-Exp1 & 16 & 1 & 4 & 3 & 4 & 4 & 1.0 to 9.8 & 0 & 0 \\
PT-Exp2 & 50 & 0 & 8 & 10 & 16 & 16 & 67.4 to 134.8 & 50 & 0 \\
PT-Exp3 & 1 & 1 & 0 & 0 & 0 & 0 & -- & 0 & 0 \\
PT-Exp4 & 14 & 1 & 4 & 2 & 4 & 3 & 1.2 to 8.9 & 2 & 12 \\
PP-Exp1 & 5 & 1 & 1 & 1 & 1 & 1 & 2.0 to 5.9 & 2 & 0 \\
PP-Exp2 & 4 & 0 & 1 & 1 & 1 & 1 & 4.8 to 13.9 & 2 & 0 \\
PP-Exp3 & 1 & 1 & 0 & 0 & 0 & 0 & -- & 0 & 0 \\
PP-Exp4 & 4 & 1 & 1 & 1 & 0 & 1 & 1.5 to 1.8 & 0 & 0 \\
ST-Exp1 & 2 & 2 & 0 & 0 & 0 & 0 & -- & 0 & 0 \\
ST-Exp2 & 2 & 2 & 0 & 0 & 0 & 0 & -- & 0 & 0 \\
ST-Exp3 & 2 & 0 & 0 & 0 & 1 & 1 & 16.6 to 25.2 & 2 & 0 \\
ST-Exp4 & 2 & 1 & 0 & 1 & 0 & 0 & 0.9 & 0 & 0 \\
SP-Exp1 & 5 & 1 & 1 & 1 & 1 & 1 & 2.2 to 3.6 & 0 & 0 \\
SP-Exp2 & 5 & 1 & 1 & 1 & 1 & 1 & 2.8 to 7.1 & 2 & 0 \\
SP-Exp3 & 6 & 1 & 2 & 1 & 1 & 1 & 0.2 to 0.5 & 0 & 0 \\
SP-Exp4 & 32 & 0 & 12 & 8 & 5 & 7 & 15.4 to 119.1 & 32 & 0 \\
\midrule
All & 151 & 14 & 35 & 30 & 35 & 37 & 0.2 to 134.8 & 92 & 12 \\
\bottomrule
\end{tabular}
\end{table}


\section{The lumped and local thresholds and the far face under step 11's draws}
\label{sec:s-step11}

Step 11 searched the lumped and the local threshold with the 21.4~s pulse at the central contact in four tests, PT-Exp1, PP-Exp3, ST-Exp2 and SP-Exp3, under the posterior mean and the nominal properties and under its eight draws (Section~\ref{sec:s-key}), each search starting from step 8's bracket, which is the posterior mean's in Table~\ref{tab:s-step11}, on a grid of 0.0125 of the test's ejected energy, coarser than the main text's searches (Table~\ref{tab:s-step11}). With no delivered heat it ran the five low-maximum tests that did not propagate under its four thermal draws at the central contact, and three of them, PP-Exp3, ST-Exp1 and ST-Exp2, at the high contact under the posterior mean and the nominal properties; Table~\ref{tab:s-farface} gives these beside the earlier steps' runs of the five under the posterior mean at the central and the low contact and of PT-Exp1 and PT-Exp3 at the high contact (steps 7 and 9).

\IfFileExists{tables/s12_step11.tex}{\begin{table}[!htbp]
\centering
\scriptsize
\caption{The lumped and the hottest point's thresholds of the four tests that step~11 searched, with the 21.4~s pulse at the central contact, under the posterior mean and the nominal properties and under the step's eight draws (Section~\ref{sec:s-key}): the bracket between the largest energy run that does not meet the criterion and the smallest that does, in kJ, on a grid of 0.0125 of the test's ejected energy (the last row); 0, met with no delivery. The main text quotes a bracket's midpoint. The kinetic draws are ranked, earliest first, by the time of the hot spot's onset in step~3's case; the thermal draws are the extremes among thirty over the roll's published ranges, each named by the property in which it is extreme and setting all three: the conductivity across and along the layers and the specific heat.}
\label{tab:s-step11}
\setlength{\tabcolsep}{4pt}
\begin{tabular}{>{\raggedright\arraybackslash}p{0.33\textwidth}rrrr}
\toprule
Draw & PT-Exp1 & PP-Exp3 & ST-Exp2 & SP-Exp3 \\
\midrule
\multicolumn{5}{l}{\textit{The lumped threshold}} \\
Posterior mean, nominal properties & (12.54, 12.89] & (1.70, 3.41] & (1.80, 3.61] & 0 \\
\addlinespace[2pt]
Kinetic draw ranked 5 of 200 & (12.54, 12.89] & (1.70, 3.41] & (0.00, 1.80] & 0 \\
Kinetic draw ranked 50 of 200 & (12.54, 12.89] & (1.70, 3.41] & (0.00, 1.80] & 0 \\
Kinetic draw ranked 151 of 200 & (12.89, 13.24] & (3.41, 5.11] & (1.80, 3.61] & 0 \\
Kinetic draw ranked 196 of 200 & (12.89, 13.24] & (3.41, 5.11] & (1.80, 3.61] & 0 \\
\addlinespace[2pt]
Least conductivity across the layers: 1.008 and 25.35 W\,m$^{-1}$\,K$^{-1}$, 1030 J\,kg$^{-1}$\,K$^{-1}$ & (12.89, 13.24] & (1.70, 3.41] & (1.80, 3.61] & 0 \\
Largest conductivity across the layers: 1.375 and 25.31 W\,m$^{-1}$\,K$^{-1}$, 1050 J\,kg$^{-1}$\,K$^{-1}$ & (12.89, 13.24] & (1.70, 3.41] & (1.80, 3.61] & 0 \\
Least specific heat: 1.177 and 25.68 W\,m$^{-1}$\,K$^{-1}$, 849 J\,kg$^{-1}$\,K$^{-1}$ & (10.80, 11.15] & (0.00, 1.70] & 0 & 0 \\
Largest specific heat: 1.196 and 25.42 W\,m$^{-1}$\,K$^{-1}$, 1109 J\,kg$^{-1}$\,K$^{-1}$ & (13.59, 13.94] & (3.41, 5.11] & (1.80, 3.61] & (0.00, 2.11] \\
\midrule
\multicolumn{5}{l}{\textit{The hottest point's threshold}} \\
Posterior mean, nominal properties & (4.18, 4.53] & (0.00, 1.70] & 0 & 0 \\
\addlinespace[2pt]
Kinetic draw ranked 5 of 200 & (4.18, 4.53] & (0.00, 1.70] & 0 & 0 \\
Kinetic draw ranked 50 of 200 & (4.18, 4.53] & (0.00, 1.70] & 0 & 0 \\
Kinetic draw ranked 151 of 200 & (4.18, 4.53] & (0.00, 1.70] & 0 & 0 \\
Kinetic draw ranked 196 of 200 & (4.18, 4.53] & (0.00, 1.70] & 0 & 0 \\
\addlinespace[2pt]
Least conductivity across the layers: 1.008 and 25.35 W\,m$^{-1}$\,K$^{-1}$, 1030 J\,kg$^{-1}$\,K$^{-1}$ & (3.83, 4.18] & 0 & 0 & 0 \\
Largest conductivity across the layers: 1.375 and 25.31 W\,m$^{-1}$\,K$^{-1}$, 1050 J\,kg$^{-1}$\,K$^{-1}$ & (4.18, 4.53] & (0.00, 1.70] & 0 & 0 \\
Least specific heat: 1.177 and 25.68 W\,m$^{-1}$\,K$^{-1}$, 849 J\,kg$^{-1}$\,K$^{-1}$ & (3.83, 4.18] & 0 & 0 & 0 \\
Largest specific heat: 1.196 and 25.42 W\,m$^{-1}$\,K$^{-1}$, 1109 J\,kg$^{-1}$\,K$^{-1}$ & (4.18, 4.53] & (0.00, 1.70] & 0 & 0 \\
\midrule
Ejected energy (kJ) & 27.88 & 136.31 & 144.26 & 168.63 \\
\bottomrule
\end{tabular}
\end{table}
}{}

\IfFileExists{tables/s13_farface.tex}{\begin{table}[!htbp]
\centering
\footnotesize
\caption{The far face's peak with no delivered heat over 900~s (\textdegree C) in the five tests that did not propagate and report low maxima, with each test's fitted trigger, against the maximum Fransson et al.\ measured: at the central contact under the posterior mean and the nominal properties and under step~11's four thermal draws, each with its three properties as in Table~\ref{tab:s-step11}, and at the low and high contacts under the posterior mean and the nominal properties.}
\label{tab:s-farface}
\setlength{\tabcolsep}{3pt}
\begin{tabular}{>{\raggedright\arraybackslash}p{0.36\textwidth}rrrrr}
\toprule
 & PT-Exp1 & PT-Exp3 & PP-Exp3 & ST-Exp1 & ST-Exp2 \\
\midrule
Measured maximum & 71 & 118 & 105 & 128 & 133 \\
\midrule
Central contact: posterior mean, nominal properties & 60.2 & 131.8 & 161.3 & 172.9 & 163.4 \\
Central contact: least conductivity across the layers: 1.008 and 25.35 W\,m$^{-1}$\,K$^{-1}$, 1030 J\,kg$^{-1}$\,K$^{-1}$ & 59.3 & 129.0 & 157.6 & 168.5 & 159.4 \\
Central contact: largest conductivity across the layers: 1.375 and 25.31 W\,m$^{-1}$\,K$^{-1}$, 1050 J\,kg$^{-1}$\,K$^{-1}$ & 60.0 & 131.2 & 160.5 & 172.1 & 162.7 \\
Central contact: least specific heat: 1.177 and 25.68 W\,m$^{-1}$\,K$^{-1}$, 849 J\,kg$^{-1}$\,K$^{-1}$ & 62.6 & 139.3 & 171.7 & 185.3 & 174.2 \\
Central contact: largest specific heat: 1.196 and 25.42 W\,m$^{-1}$\,K$^{-1}$, 1109 J\,kg$^{-1}$\,K$^{-1}$ & 58.8 & 127.5 & 155.5 & 166.3 & 157.4 \\
\midrule
Low contact, posterior mean and nominal properties & 50.1 & 100.4 & 121.6 & 127.0 & 120.2 \\
High contact, posterior mean and nominal properties & 61.9 & 137.9 & 169.9 & 183.7 & 173.0 \\
\bottomrule
\end{tabular}
\end{table}
}{}

\section{The archived runs one by one}
\label{sec:s-archived-runs}

Table~\ref{tab:s-archived-runs} lists the 151 runs of Section~\ref{sec:s-archived}, by test, delivery and energy, with the step whose search made each and both models' outcomes. Where a model's resolved neighbour runs away, the table gives the time its mean reaches 255.7\degC; where it does not, the mean's peak by the run's end.

\IfFileExists{tables/s14_archived_runs.tex}{{\scriptsize\setlength{\tabcolsep}{4pt}
\begin{longtable}{llrrrll}
\caption{The 151 archived three-dimensional runs of Table~\ref{tab:s-archived} one by one, at the central contact: the test, the delivery (none, the 21.4~s pulse, or constant power over the time given), the energy delivered to the facing half, the step whose search made the run, how far it was carried, and each model's outcome, the time at which the resolved neighbour's mean reaches 255.7\textdegree C or, where it does not, the mean's peak. A dagger marks the six runs that repeat another's configuration in a different step; they are counted as they were run. Two energies differ in the last digit from the needs of Tables~\ref{tab:s-thresholds} and~\ref{tab:s-needs}: PT-Exp1's over 360~s is the same run, whose energy the needs file records to 0.1~J, and SP-Exp3's with the pulse is step 9's run at the fraction its search interpolated, while the need is that of step 12's run, which met the measured time.}\label{tab:s-archived-runs}\\
\toprule
Test & Delivery & Energy (kJ) & Step & To (s) & Three-dimensional & Cross-section \\
\midrule
\endfirsthead
\multicolumn{7}{l}{\textit{Continued}} \\
\toprule
Test & Delivery & Energy (kJ) & Step & To (s) & Three-dimensional & Cross-section \\
\midrule
\endhead
PT-Exp1 & none & 0 & 10 & 900 & peak 61.5\textdegree C & peak 60.3\textdegree C \\
PT-Exp1 & pulse & 1.05 & 10 & 900 & peak 66.7\textdegree C & peak 65.7\textdegree C \\
PT-Exp1 & pulse & 2.09 & 9 & 900 & peak 73.2\textdegree C & peak 72.5\textdegree C \\
PT-Exp1$^\dagger$ & pulse & 2.09 & 10 & 900 & peak 73.2\textdegree C & peak 72.5\textdegree C \\
PT-Exp1 & pulse & 4.18 & 10 & 900 & peak 90.6\textdegree C & peak 90.5\textdegree C \\
PT-Exp1 & 120~s & 1.52 & 9 & 900 & peak 71.4\textdegree C & peak 70.3\textdegree C \\
PT-Exp1 & 120~s & 3.05 & 9 & 900 & peak 85.8\textdegree C & peak 84.7\textdegree C \\
PT-Exp1 & 120~s & 6.10 & 9 & 900 & peak 121.5\textdegree C & peak 121.3\textdegree C \\
PT-Exp1 & 240~s & 1.01 & 9 & 900 & peak 69.4\textdegree C & peak 68.0\textdegree C \\
PT-Exp1 & 240~s & 2.01 & 9 & 900 & peak 79.2\textdegree C & peak 78.1\textdegree C \\
PT-Exp1 & 240~s & 4.03 & 9 & 900 & peak 98.8\textdegree C & peak 98.6\textdegree C \\
PT-Exp1 & 240~s & 8.06 & 9 & 900 & peak 137.4\textdegree C & peak 139.0\textdegree C \\
PT-Exp1 & 360~s & 1.22 & 9 & 900 & peak 71.6\textdegree C & peak 70.9\textdegree C \\
PT-Exp1 & 360~s & 2.45 & 9 & 900 & peak 81.7\textdegree C & peak 81.6\textdegree C \\
PT-Exp1 & 360~s & 4.90 & 9 & 900 & peak 101.6\textdegree C & peak 102.9\textdegree C \\
PT-Exp1 & 360~s & 9.80 & 9 & 900 & peak 140.6\textdegree C & peak 144.4\textdegree C \\
PT-Exp2 & pulse & 67.39 & 9 & 900 & runs away, 6.5~s & runs away, 6.5~s \\
PT-Exp2 & pulse & 71.60 & 9 & 900 & runs away, 6.3~s & runs away, 6.3~s \\
PT-Exp2 & pulse & 75.81 & 9 & 900 & runs away, 6.1~s & runs away, 6.1~s \\
PT-Exp2 & pulse & 80.03 & 9 & 900 & runs away, 6.0~s & runs away, 6.0~s \\
PT-Exp2 & pulse & 84.24 & 9 & 900 & runs away, 5.8~s & runs away, 5.8~s \\
PT-Exp2 & pulse & 92.66 & 9 & 900 & runs away, 5.6~s & runs away, 5.6~s \\
PT-Exp2 & pulse & 101.09 & 9 & 900 & runs away, 5.3~s & runs away, 5.3~s \\
PT-Exp2 & pulse & 117.93 & 9 & 900 & runs away, 5.0~s & runs away, 5.0~s \\
PT-Exp2 & 120~s & 67.39 & 9 & 900 & runs away, 20.1~s & runs away, 20.0~s \\
PT-Exp2 & 120~s & 75.81 & 9 & 900 & runs away, 18.2~s & runs away, 18.0~s \\
PT-Exp2 & 120~s & 84.24 & 9 & 900 & runs away, 16.5~s & runs away, 16.4~s \\
PT-Exp2 & 120~s & 88.45 & 9 & 900 & runs away, 15.8~s & runs away, 15.6~s \\
PT-Exp2 & 120~s & 92.66 & 9 & 900 & runs away, 15.1~s & runs away, 15.0~s \\
PT-Exp2 & 120~s & 94.77 & 9 & 900 & runs away, 14.8~s & runs away, 14.7~s \\
PT-Exp2 & 120~s & 96.87 & 9 & 900 & runs away, 14.5~s & runs away, 14.4~s \\
PT-Exp2 & 120~s & 98.98 & 9 & 900 & runs away, 14.2~s & runs away, 14.1~s \\
PT-Exp2 & 120~s & 101.09 & 9 & 900 & runs away, 13.9~s & runs away, 13.8~s \\
PT-Exp2 & 120~s & 117.93 & 9 & 900 & runs away, 12.1~s & runs away, 11.9~s \\
PT-Exp2 & 240~s & 67.39 & 9 & 900 & runs away, 37.2~s & runs away, 37.2~s \\
PT-Exp2 & 240~s & 84.24 & 9 & 900 & runs away, 30.5~s & runs away, 30.4~s \\
PT-Exp2 & 240~s & 101.09 & 9 & 900 & runs away, 25.9~s & runs away, 25.8~s \\
PT-Exp2 & 240~s & 117.93 & 9 & 900 & runs away, 22.6~s & runs away, 22.5~s \\
PT-Exp2 & 240~s & 122.15 & 9 & 900 & runs away, 22.0~s & runs away, 21.9~s \\
PT-Exp2 & 240~s & 126.36 & 9 & 900 & runs away, 21.3~s & runs away, 21.2~s \\
PT-Exp2 & 240~s & 128.46 & 9 & 900 & runs away, 21.0~s & runs away, 20.9~s \\
PT-Exp2 & 240~s & 130.57 & 9 & 900 & runs away, 20.7~s & runs away, 20.6~s \\
PT-Exp2 & 240~s & 131.62 & 9 & 900 & runs away, 20.6~s & runs away, 20.5~s \\
PT-Exp2 & 240~s & 132.68 & 9 & 900 & runs away, 20.4~s & runs away, 20.3~s \\
PT-Exp2 & 240~s & 133.20 & 9 & 900 & runs away, 20.4~s & runs away, 20.2~s \\
PT-Exp2 & 240~s & 133.73 & 9 & 900 & runs away, 20.3~s & runs away, 20.2~s \\
PT-Exp2 & 240~s & 133.99 & 9 & 900 & runs away, 20.3~s & runs away, 20.1~s \\
PT-Exp2 & 240~s & 134.26 & 9 & 900 & runs away, 20.2~s & runs away, 20.1~s \\
PT-Exp2 & 240~s & 134.52 & 9 & 900 & runs away, 20.2~s & runs away, 20.1~s \\
PT-Exp2 & 240~s & 134.78 & 9 & 900 & runs away, 20.1~s & runs away, 20.0~s \\
PT-Exp2 & 360~s & 67.39 & 9 & 900 & runs away, 52.7~s & runs away, 53.0~s \\
PT-Exp2 & 360~s & 84.24 & 9 & 900 & runs away, 43.6~s & runs away, 43.7~s \\
PT-Exp2 & 360~s & 101.09 & 9 & 900 & runs away, 37.2~s & runs away, 37.2~s \\
PT-Exp2 & 360~s & 109.51 & 9 & 900 & runs away, 34.6~s & runs away, 34.6~s \\
PT-Exp2 & 360~s & 117.93 & 9 & 900 & runs away, 32.4~s & runs away, 32.4~s \\
PT-Exp2 & 360~s & 122.15 & 9 & 900 & runs away, 31.4~s & runs away, 31.4~s \\
PT-Exp2 & 360~s & 126.36 & 9 & 900 & runs away, 30.5~s & runs away, 30.4~s \\
PT-Exp2 & 360~s & 128.46 & 9 & 900 & runs away, 30.0~s & runs away, 30.0~s \\
PT-Exp2 & 360~s & 130.57 & 9 & 900 & runs away, 29.6~s & runs away, 29.5~s \\
PT-Exp2 & 360~s & 131.62 & 9 & 900 & runs away, 29.4~s & runs away, 29.3~s \\
PT-Exp2 & 360~s & 132.68 & 9 & 900 & runs away, 29.2~s & runs away, 29.1~s \\
PT-Exp2 & 360~s & 133.73 & 9 & 900 & runs away, 29.0~s & runs away, 28.9~s \\
PT-Exp2 & 360~s & 133.99 & 9 & 900 & runs away, 28.9~s & runs away, 28.9~s \\
PT-Exp2 & 360~s & 134.26 & 9 & 900 & runs away, 28.9~s & runs away, 28.8~s \\
PT-Exp2 & 360~s & 134.52 & 9 & 900 & runs away, 28.8~s & runs away, 28.8~s \\
PT-Exp2 & 360~s & 134.78 & 9 & 900 & runs away, 28.8~s & runs away, 28.7~s \\
PT-Exp3 & none & 0 & 10 & 900 & peak 136.3\textdegree C & peak 132.6\textdegree C \\
PT-Exp4 & none & 0 & 10 & 332 & peak 158.1\textdegree C & peak 153.7\textdegree C \\
PT-Exp4 & pulse & 1.16 & 10 & 332 & peak 164.3\textdegree C & peak 160.1\textdegree C \\
PT-Exp4 & pulse & 1.55 & 10 & 332 & peak 166.4\textdegree C & peak 162.4\textdegree C \\
PT-Exp4 & pulse & 4.66 & 10 & 332 & peak 187.0\textdegree C & peak 184.1\textdegree C \\
PT-Exp4 & pulse & 5.04 & 10 & 332 & peak 190.6\textdegree C & peak 187.9\textdegree C \\
PT-Exp4 & 120~s & 3.49 & 10 & 332 & peak 182.7\textdegree C & peak 178.6\textdegree C \\
PT-Exp4 & 120~s & 3.88 & 10 & 332 & peak 185.9\textdegree C & peak 181.8\textdegree C \\
PT-Exp4 & 240~s & 3.49 & 9 & 900 & peak 189.8\textdegree C & peak 185.3\textdegree C \\
PT-Exp4$^\dagger$ & 240~s & 3.49 & 10 & 332 & peak 189.8\textdegree C & peak 185.3\textdegree C \\
PT-Exp4 & 240~s & 6.21 & 10 & 332 & peak 241.2\textdegree C & peak 220.5\textdegree C \\
PT-Exp4 & 240~s & 6.60 & 10 & 332 & runs away, 263.2~s & runs away, 293.3~s \\
PT-Exp4 & 360~s & 5.04 & 9 & 900 & peak 196.8\textdegree C & peak 196.5\textdegree C \\
PT-Exp4$^\dagger$ & 360~s & 5.04 & 10 & 332 & peak 196.3\textdegree C & peak 195.0\textdegree C \\
PT-Exp4 & 360~s & 8.93 & 10 & 332 & runs away, 283.2~s & runs away, 293.1~s \\
PP-Exp1 & none & 0 & 10 & 900 & peak 183.6\textdegree C & peak 178.3\textdegree C \\
PP-Exp1 & pulse & 2.97 & 9 & 900 & peak 204.3\textdegree C & peak 198.8\textdegree C \\
PP-Exp1 & 120~s & 1.98 & 9 & 900 & peak 199.3\textdegree C & peak 193.5\textdegree C \\
PP-Exp1 & 240~s & 3.97 & 9 & 900 & runs away, 243.5~s & runs away, 278.0~s \\
PP-Exp1 & 360~s & 5.95 & 9 & 900 & runs away, 243.1~s & runs away, 266.2~s \\
PP-Exp2 & pulse & 5.96 & 9 & 900 & peak 204.5\textdegree C & peak 202.2\textdegree C \\
PP-Exp2 & 120~s & 4.77 & 9 & 900 & peak 197.1\textdegree C & peak 192.6\textdegree C \\
PP-Exp2 & 240~s & 7.55 & 9 & 900 & runs away, 219.6~s & runs away, 229.4~s \\
PP-Exp2 & 360~s & 13.91 & 9 & 900 & runs away, 186.9~s & runs away, 194.7~s \\
PP-Exp3 & none & 0 & 10 & 900 & peak 167.4\textdegree C & peak 162.7\textdegree C \\
PP-Exp4 & none & 0 & 10 & 900 & peak 174.2\textdegree C & peak 169.3\textdegree C \\
PP-Exp4 & pulse & 1.82 & 9 & 900 & peak 184.7\textdegree C & peak 180.2\textdegree C \\
PP-Exp4 & 120~s & 1.46 & 9 & 900 & peak 184.3\textdegree C & peak 179.5\textdegree C \\
PP-Exp4 & 360~s & 1.82 & 9 & 900 & peak 187.2\textdegree C & peak 183.7\textdegree C \\
ST-Exp1 & none & 0 & 9 & 900 & peak 179.5\textdegree C & peak 179.1\textdegree C \\
ST-Exp1$^\dagger$ & none & 0 & 10 & 900 & peak 179.5\textdegree C & peak 179.1\textdegree C \\
ST-Exp2 & none & 0 & 9 & 900 & peak 169.5\textdegree C & peak 169.1\textdegree C \\
ST-Exp2$^\dagger$ & none & 0 & 10 & 900 & peak 169.5\textdegree C & peak 169.1\textdegree C \\
ST-Exp3 & 240~s & 16.63 & 9 & 900 & runs away, 207.9~s & runs away, 202.4~s \\
ST-Exp3 & 360~s & 25.19 & 9 & 900 & runs away, 205.5~s & runs away, 200.2~s \\
ST-Exp4 & none & 0 & 10 & 900 & peak 174.6\textdegree C & peak 174.2\textdegree C \\
ST-Exp4 & 120~s & 0.95 & 9 & 900 & peak 181.1\textdegree C & peak 180.8\textdegree C \\
SP-Exp1 & none & 0 & 10 & 900 & peak 170.3\textdegree C & peak 169.9\textdegree C \\
SP-Exp1 & pulse & 2.73 & 9 & 900 & peak 186.5\textdegree C & peak 186.8\textdegree C \\
SP-Exp1 & 120~s & 2.17 & 9 & 900 & peak 185.7\textdegree C & peak 185.5\textdegree C \\
SP-Exp1 & 240~s & 2.27 & 9 & 900 & peak 191.1\textdegree C & peak 190.4\textdegree C \\
SP-Exp1 & 360~s & 3.63 & 9 & 900 & peak 198.5\textdegree C & peak 201.8\textdegree C \\
SP-Exp2 & none & 0 & 10 & 900 & peak 174.4\textdegree C & peak 174.1\textdegree C \\
SP-Exp2 & pulse & 3.78 & 9 & 900 & peak 199.5\textdegree C & peak 200.5\textdegree C \\
SP-Exp2 & 120~s & 2.84 & 9 & 900 & peak 196.1\textdegree C & peak 196.2\textdegree C \\
SP-Exp2 & 240~s & 5.68 & 9 & 900 & runs away, 221.8~s & runs away, 223.1~s \\
SP-Exp2 & 360~s & 7.09 & 9 & 900 & runs away, 253.9~s & runs away, 253.3~s \\
SP-Exp3 & none & 0 & 10 & 900 & peak 186.2\textdegree C & peak 185.9\textdegree C \\
SP-Exp3 & pulse & 0.27 & 9 & 900 & peak 187.8\textdegree C & peak 187.6\textdegree C \\
SP-Exp3$^\dagger$ & pulse & 0.27 & 11 & 900 & peak 187.8\textdegree C & peak 187.6\textdegree C \\
SP-Exp3 & 120~s & 0.22 & 9 & 900 & peak 187.8\textdegree C & peak 187.5\textdegree C \\
SP-Exp3 & 240~s & 0.26 & 9 & 900 & peak 188.6\textdegree C & peak 188.3\textdegree C \\
SP-Exp3 & 360~s & 0.53 & 9 & 900 & peak 190.1\textdegree C & peak 190.1\textdegree C \\
SP-Exp4 & pulse & 15.37 & 9 & 900 & runs away, 14.0~s & runs away, 13.8~s \\
SP-Exp4 & pulse & 30.75 & 9 & 900 & runs away, 9.5~s & runs away, 9.5~s \\
SP-Exp4 & pulse & 31.71 & 9 & 900 & runs away, 9.3~s & runs away, 9.3~s \\
SP-Exp4 & pulse & 32.67 & 9 & 900 & runs away, 9.2~s & runs away, 9.2~s \\
SP-Exp4 & pulse & 33.63 & 9 & 900 & runs away, 9.1~s & runs away, 9.1~s \\
SP-Exp4 & pulse & 34.59 & 9 & 900 & runs away, 9.0~s & runs away, 8.9~s \\
SP-Exp4 & pulse & 36.51 & 9 & 900 & runs away, 8.7~s & runs away, 8.7~s \\
SP-Exp4 & pulse & 38.43 & 9 & 900 & runs away, 8.5~s & runs away, 8.5~s \\
SP-Exp4 & pulse & 42.27 & 9 & 900 & runs away, 8.1~s & runs away, 8.1~s \\
SP-Exp4 & pulse & 46.12 & 9 & 900 & runs away, 7.8~s & runs away, 7.8~s \\
SP-Exp4 & pulse & 53.80 & 9 & 900 & runs away, 7.2~s & runs away, 7.2~s \\
SP-Exp4 & pulse & 61.49 & 9 & 900 & runs away, 6.8~s & runs away, 6.8~s \\
SP-Exp4 & 120~s & 30.75 & 9 & 900 & runs away, 41.8~s & runs away, 41.4~s \\
SP-Exp4 & 120~s & 38.43 & 9 & 900 & runs away, 34.1~s & runs away, 33.8~s \\
SP-Exp4 & 120~s & 40.35 & 9 & 900 & runs away, 32.6~s & runs away, 32.3~s \\
SP-Exp4 & 120~s & 42.27 & 9 & 900 & runs away, 31.2~s & runs away, 30.9~s \\
SP-Exp4 & 120~s & 44.20 & 9 & 900 & runs away, 29.9~s & runs away, 29.7~s \\
SP-Exp4 & 120~s & 46.12 & 9 & 900 & runs away, 28.8~s & runs away, 28.5~s \\
SP-Exp4 & 120~s & 53.80 & 9 & 900 & runs away, 25.0~s & runs away, 24.8~s \\
SP-Exp4 & 120~s & 61.49 & 9 & 900 & runs away, 22.1~s & runs away, 22.0~s \\
SP-Exp4 & 240~s & 61.49 & 9 & 900 & runs away, 41.8~s & runs away, 41.4~s \\
SP-Exp4 & 240~s & 69.18 & 9 & 900 & runs away, 37.5~s & runs away, 37.2~s \\
SP-Exp4 & 240~s & 76.86 & 9 & 900 & runs away, 34.1~s & runs away, 33.8~s \\
SP-Exp4 & 240~s & 92.24 & 9 & 900 & runs away, 28.8~s & runs away, 28.5~s \\
SP-Exp4 & 240~s & 107.61 & 9 & 900 & runs away, 25.0~s & runs away, 24.8~s \\
SP-Exp4 & 360~s & 61.49 & 9 & 900 & runs away, 59.9~s & runs away, 59.6~s \\
SP-Exp4 & 360~s & 92.24 & 9 & 900 & runs away, 41.8~s & runs away, 41.4~s \\
SP-Exp4 & 360~s & 99.92 & 9 & 900 & runs away, 38.8~s & runs away, 38.5~s \\
SP-Exp4 & 360~s & 107.61 & 9 & 900 & runs away, 36.3~s & runs away, 36.0~s \\
SP-Exp4 & 360~s & 111.45 & 9 & 900 & runs away, 35.2~s & runs away, 34.8~s \\
SP-Exp4 & 360~s & 115.30 & 9 & 900 & runs away, 34.1~s & runs away, 33.8~s \\
SP-Exp4 & 360~s & 119.14 & 9 & 900 & runs away, 33.1~s & runs away, 32.8~s \\
\bottomrule
\end{longtable}
}
}{}

\section{The resolved neighbour with no delivered heat}
\label{sec:s-nodelivery}

Table~\ref{tab:s-nodelivery} gives, for every test at the three contacts, the resolved neighbour's outcome with no delivered heat under the test's fitted trigger, and how it is established, which Section~3.1 of the main text summarises. No test's resolved neighbour runs away within 900~s at any contact. At the low contact the hottest point stays below 255.7\degC{} in step 7's three-dimensional runs to 900~s in all sixteen tests, so the mean cannot reach it; at the central contact it does so in eight, and the other eight were carried to 900~s in steps 9 and 10; at the high contact it does so in three, step 11 carried three more to 900~s, and the other ten rest on the reactive cross-section alone, which was set against the three-dimensional model only at the central contact.

\IfFileExists{tables/s15_nodelivery.tex}{\begin{table}[!htbp]
\centering
\footnotesize
\caption{The resolved neighbour with no delivered heat, test by test at the three contacts, with each test's fitted trigger under the posterior mean and the nominal properties: the peak of its mean temperature within 900~s (\textdegree C) and how its outcome, no runaway in every case, is established. $^{h}$, the hottest point stays below 255.7\textdegree C in a three-dimensional run to 900~s, so the mean cannot reach it; $^{r}$, a three-dimensional run carried to 900~s; $^{x}$, the reactive cross-section alone (Section~3.7 of the main text). The peak is the three-dimensional run's where there is one and the cross-section's otherwise.}
\label{tab:s-nodelivery}
\setlength{\tabcolsep}{6pt}
\begin{tabular}{lrrr}
\toprule
Test & Low contact & Central contact & High contact \\
\midrule
PT-Exp1 & 51.0$^{h}$ & 61.5$^{h}$ & 63.2$^{h}$ \\
PT-Exp2 & 125.0$^{h}$ & 166.3$^{h}$ & 171.4$^{x}$ \\
PT-Exp3 & 103.6$^{h}$ & 136.3$^{h}$ & 142.7$^{h}$ \\
PT-Exp4 & 119.1$^{h}$ & 158.1$^{h}$ & 162.7$^{x}$ \\
PP-Exp1 & 137.5$^{h}$ & 183.6$^{r}$ & 190.1$^{x}$ \\
PP-Exp2 & 120.8$^{h}$ & 160.5$^{h}$ & 165.2$^{x}$ \\
PP-Exp3 & 125.8$^{h}$ & 167.4$^{h}$ & 176.5$^{r}$ \\
PP-Exp4 & 130.7$^{h}$ & 174.2$^{r}$ & 179.9$^{x}$ \\
ST-Exp1 & 131.2$^{h}$ & 179.5$^{r}$ & 190.8$^{r}$ \\
ST-Exp2 & 124.1$^{h}$ & 169.5$^{r}$ & 179.5$^{r}$ \\
ST-Exp3 & 68.0$^{h}$ & 88.2$^{h}$ & 91.7$^{h}$ \\
ST-Exp4 & 127.7$^{h}$ & 174.6$^{r}$ & 185.2$^{x}$ \\
SP-Exp1 & 124.7$^{h}$ & 170.3$^{r}$ & 180.4$^{x}$ \\
SP-Exp2 & 127.6$^{h}$ & 174.4$^{r}$ & 185.0$^{x}$ \\
SP-Exp3 & 136.0$^{h}$ & 186.2$^{r}$ & 199.2$^{x}$ \\
SP-Exp4 & 112.9$^{h}$ & 153.8$^{h}$ & 162.3$^{x}$ \\
\bottomrule
\end{tabular}
\end{table}
}{}

\section{The three-dimensional needs and PT-Exp1's threshold}
\label{sec:s-step16}

Step 16 searched in three dimensions the needs with the pulse that the main text had taken from the reactive cross-section, PT-Exp1's threshold with the pulse, and, as computing allowed, the needs with constant power, with the configuration, kinetics, properties, meshes, time steps and rules of the runs they extend; its predictions were pushed before any of its runs (Table~\ref{tab:s-steps}). A run is early when the resolved neighbour's mean reaches 255.7\degC{} no later than 10~s after the measured time, and late when it reaches it later or not within 900~s; a need's bracket runs from the highest late energy below every early one to the lowest early energy. Every three-dimensional run already made in the same configuration counted, the archived runs of Table~\ref{tab:s-archived} and step 15's threshold searches. Each next run lay on a grid from the cross-section's need in steps of 0.25\,\% of the test's ejected energy, upward while no run was early and downward while none was late, and then halved the bracket until it was at most 0.125\,\% of the ejected energy wide. PT-Exp1's threshold was searched as step 15 searched those of Table~3 of the main text, opening 5\,\% either side of the cross-section's value with steps of 10\,\% of it. One need, the one with the latest measured time among those whose bracket's lower end also runs away, was refined: both ends of its bracket were run again with the mesh refined by a factor of 1.5 and every step limit halved together. The needs with constant power were searched by the same rule, over 240, 120 and 360~s in that order and test by test, until no search was left or the runs made had taken six hours of a worker's time. Table~\ref{tab:s-step16} gives every run.

Before V and X, step 16 made 33 runs, T's two and N's 31, with 18 archived runs counted beside N's; every bracket's two ends are runs of step 16 (Table~\ref{tab:s-step16}). Table~3 of the main text gives the eight needs with the pulse beside the cross-section's and the lumped needs. Against the bracket's midpoint, the cross-section's need, the upper end of its own search's bracket, 0.5\,\% of the test's ejected energy wide, lies 1.5\,\% below to 8.6\,\% above it (P1, held); it lies above it in seven of the eight, and in SP-Exp1 below it (P2, failed), the run at the cross-section's need, 5.96~kJ, running away at 250.6~s against the measured 223~s. Every bracket lies within 4.45 to 7.54~kJ (P3, held), 1.46 to 4.33~kJ above the lumped need (P4 and P5, held). In PT-Exp4 and ST-Exp4 the bracket's lower end does not run away within 900~s, the mean peaking at 243.1 and 248.6\degC, and its upper end runs away at 250.8 and 254.6~s, before the measured 272 and 342~s, so that no run reproduces the measured time and the same two runs bracket the threshold: PT-Exp4's inside its threshold search's bracket, (6.71, 7.42], and ST-Exp4's, which no threshold search reached, between 5.32 and 5.51~kJ; this reading was made after the runs, and no prediction covered it. PT-Exp1's threshold closed at the first two runs (P6 and P7, held): at 10.24~kJ the mean peaks at 218.2\degC, and at 11.32~kJ it reaches 255.7\degC{} at 45.8~s while the lumped neighbour peaks at 204.2\degC; the bracket's midpoint is therefore the cross-section's value, and the three-dimensional runs say only that the threshold lies within 5\,\% of it. V refined PP-Exp4's need, the one with the latest measured time among those whose bracket's lower end runs away: with the mesh refined by a factor of 1.5 and every step limit halved together, the runaway at the bracket's lower end came at 306.8~s against 298.8~s, and at its upper end at 252.1~s against 249.3~s, later by 2.7 and 1.1\,\%, both ends keeping their class, so that the need's bracket is unchanged (P8 and P9, held). X bracketed the needs with constant power over 240~s in all eight tests and over 120~s in six, SP-Exp2's and SP-Exp3's left unfinished and the 360~s delivery not reached when the six worker-hours were used (6.16, with the runs then under way). With the bracket counted as delivered by 10~s after the measured time and the cross-section's need by its own runaway, as registered, the cross-section's need lies 0.9 to 15.7\,\% above the bracket's midpoint, beyond 10\,\% in PP-Exp1 over 240~s (15.7\,\%) and PP-Exp4 over 120~s (14.4\,\%) (P10, failed, as its file expected), and every bracket's lower end lies above the lumped need (P11, held). Of the eleven predictions, nine held and two failed, P2 and P10.

\IfFileExists{tables/s16_step16.tex}{{\scriptsize\setlength{\tabcolsep}{4pt}
\begin{longtable}{llllll}
\caption{The runs of step 16, search by search: N, the needs with the 21.4~s pulse; T, PT-Exp1's threshold; V, the need refined (the mesh by a factor of 1.5 and every step limit halved, together); X, the needs with constant power. Each gives the energy delivered to the facing half, when the resolved neighbour's mean reaches 255.7\textdegree C (or its peak where it does not within 900~s), and, for a need, whether that comes no later than 10~s after the measured time (early) or later or not at all (late). An asterisk marks an archived run that closes a bracket. The brackets are in Table~3 of the main text and Table~\ref{tab:s-needs}.}\label{tab:s-step16}\\
\toprule
Search & Test & Delivery & Energy (kJ) & Mean reaches 255.7\textdegree C & Class \\
\midrule
\endfirsthead
\multicolumn{6}{l}{\textit{Continued}} \\
\toprule
Search & Test & Delivery & Energy (kJ) & Mean reaches 255.7\textdegree C & Class \\
\midrule
\endhead
T & PT-Exp1 & pulse & 10.242 & no; peak 218.2\textdegree C &  \\
T & PT-Exp1 & pulse & 11.320 & 45.8~s &  \\
N & PP-Exp1 & pulse & 4.505 & 264.7~s & late \\
N & PP-Exp1 & pulse & 5.270 & 184.8~s & late \\
N & PP-Exp1 & pulse & 5.462 & 172.2~s & late \\
N & PP-Exp1 & pulse & 5.653 & 160.8~s & early \\
N & PP-Exp1 & pulse & 6.036 & 140.8~s & early \\
N & PP-Exp2 & pulse & 7.423 & 160.6~s & late \\
N & PP-Exp2 & pulse & 7.536 & 149.9~s & early \\
N & PP-Exp2 & pulse & 7.649 & 140.4~s & early \\
N & PP-Exp4 & pulse & 3.934 & no; peak 199.9\textdegree C & late \\
N & PP-Exp4 & pulse & 4.991 & no; peak 213.8\textdegree C & late \\
N & PP-Exp4 & pulse & 5.520 & 298.8~s & late \\
N & PP-Exp4 & pulse & 5.652 & 249.3~s & early \\
N & PP-Exp4 & pulse & 5.784 & 224.4~s & early \\
N & PP-Exp4 & pulse & 6.048 & 191.6~s & early \\
N & PT-Exp4 & pulse & 7.062 & no; peak 243.1\textdegree C & late \\
N & PT-Exp4 & pulse & 7.151 & 250.8~s & early \\
N & PT-Exp4 & pulse & 7.239 & 216.4~s & early \\
N & SP-Exp1 & pulse & 5.965 & 250.6~s & late \\
N & SP-Exp1 & pulse & 6.146 & 214.4~s & early \\
N & SP-Exp1 & pulse & 6.328 & 191.1~s & early \\
N & SP-Exp2 & pulse & 5.139 & no; peak 223.4\textdegree C & late \\
N & SP-Exp2 & pulse & 5.817 & 207.3~s & late \\
N & SP-Exp2 & pulse & 6.156 & 174.4~s & late \\
N & SP-Exp2 & pulse & 6.326 & 161.6~s & late \\
N & SP-Exp2 & pulse & 6.495 & 150.3~s & early \\
N & SP-Exp3 & pulse & 4.319 & 234.8~s & late \\
N & SP-Exp3 & pulse & 4.451 & 221.9~s & late \\
N & SP-Exp3 & pulse & 4.582 & 210.6~s & early \\
N & ST-Exp4 & pulse & 5.129 & no; peak 223.9\textdegree C & late \\
N & ST-Exp4 & pulse & 5.319 & no; peak 248.6\textdegree C & late \\
N & ST-Exp4 & pulse & 5.508 & 254.6~s & early \\
V & PP-Exp4 & pulse & 5.520 & 306.8~s & late \\
V & PP-Exp4 & pulse & 5.652 & 252.1~s & early \\
X & PP-Exp1 & 120~s & 3.701 & 280.1~s & late \\
X & PP-Exp1 & 120~s & 4.560 & 181.0~s & late \\
X & PP-Exp1 & 120~s & 4.775 & 168.3~s & late \\
X & PP-Exp1 & 120~s & 4.883 & 162.8~s & late \\
X & PP-Exp1 & 120~s & 4.990 & 157.7~s & early \\
X & PP-Exp1 & 120~s & 5.419 & 141.1~s & early \\
X & PP-Exp1 & 240~s & 6.726 & 176.2~s & late \\
X & PP-Exp1 & 240~s & 7.415 & 166.4~s & late \\
X & PP-Exp1 & 240~s & 7.588 & 164.1~s & late \\
X & PP-Exp1 & 240~s & 7.760 & 162.0~s & early \\
X & PP-Exp1 & 240~s & 8.105 & 157.8~s & early \\
X & PP-Exp1 & 240~s & 9.485 & 143.6~s & early \\
X & PP-Exp2 & 120~s & 6.257 & no; peak 234.2\textdegree C & late \\
X & PP-Exp2 & 120~s & 7.000 & 166.2~s & late \\
X & PP-Exp2 & 120~s & 7.093 & 159.9~s & late \\
X & PP-Exp2 & 120~s & 7.186 & 154.5~s & early \\
X & PP-Exp2 & 120~s & 7.372 & 145.6~s & early \\
X & PP-Exp2 & 120~s & 7.743 & 132.2~s & early \\
X & PP-Exp2 & 240~s & 10.026 & 176.1~s & late \\
X & PP-Exp2 & 240~s & 11.264 & 161.4~s & late \\
X & PP-Exp2 & 240~s & 11.573 & 158.1~s & late \\
X & PP-Exp2 & 240~s & 11.727 & 156.5~s & early \\
X & PP-Exp2 & 240~s & 11.882 & 155.0~s & early \\
X & PP-Exp2 & 240~s & 12.501 & 149.2~s & early \\
X & PP-Exp4 & 120~s & 3.463 & no; peak 201.7\textdegree C & late \\
X & PP-Exp4 & 120~s & 4.463 & no; peak 219.9\textdegree C & late \\
X & PP-Exp4 & 120~s & 4.713 & no; peak 249.7\textdegree C & late \\
X & PP-Exp4 & 120~s & 4.838 & 269.3~s & early \\
X & PP-Exp4 & 120~s & 4.963 & 239.9~s & early \\
X & PP-Exp4 & 120~s & 5.463 & 183.1~s & early \\
X & PP-Exp4 & 240~s & 4.297 & no; peak 246.4\textdegree C & late \\
X & PP-Exp4 & 240~s & 4.479 & 292.7~s & late \\
X & PP-Exp4 & 240~s & 4.661 & 268.5~s & early \\
X & PP-Exp4 & 240~s & 5.025 & 245.5~s & early \\
X & PT-Exp4 & 120~s & 5.539 & no; peak 202.8\textdegree C & late \\
X & PT-Exp4 & 120~s & 6.368 & no; peak 220.3\textdegree C & late \\
X & PT-Exp4 & 120~s & 6.575 & no; peak 244.5\textdegree C & late \\
X & PT-Exp4 & 120~s & 6.679 & 249.3~s & early \\
X & PT-Exp4 & 120~s & 6.782 & 219.3~s & early \\
X & PT-Exp4 & 120~s & 7.197 & 167.7~s & early \\
X & PT-Exp4 & 240~s & 6.403 & 285.8~s & late \\
X & PT-Exp4 & 240~s & 6.500 & 272.3~s & early \\
X & SP-Exp1 & 120~s & 3.784 & no; peak 200.2\textdegree C & late \\
X & SP-Exp1 & 120~s & 4.592 & no; peak 211.5\textdegree C & late \\
X & SP-Exp1 & 120~s & 4.996 & no; peak 228.4\textdegree C & late \\
X & SP-Exp1 & 120~s & 5.198 & 281.7~s & late \\
X & SP-Exp1 & 120~s & 5.299 & 249.5~s & late \\
X & SP-Exp1 & 120~s & 5.400 & 229.7~s & early \\
X & SP-Exp1 & 240~s & 4.310 & no; peak 215.4\textdegree C & late \\
X & SP-Exp1 & 240~s & 5.329 & 251.2~s & late \\
X & SP-Exp1 & 240~s & 5.584 & 240.0~s & late \\
X & SP-Exp1 & 240~s & 5.711 & 235.8~s & late \\
X & SP-Exp1 & 240~s & 5.839 & 231.9~s & early \\
X & SP-Exp1 & 240~s & 6.348 & 217.8~s & early \\
X & SP-Exp2 & 120~s & 4.359 & no; peak 219.4\textdegree C & late \\
X & SP-Exp2 & 120~s & 5.119 & 208.0~s & late \\
X & SP-Exp2 & 120~s & 5.880 & 153.5~s & early \\
X & SP-Exp2 & 240~s & 7.687 & 180.5~s & late \\
X & SP-Exp2 & 240~s & 8.693 & 166.4~s & late \\
X & SP-Exp2 & 240~s & 9.196 & 160.3~s & late \\
X & SP-Exp2 & 240~s & 9.321 & 158.9~s & early \\
X & SP-Exp2 & 240~s & 9.447 & 157.5~s & early \\
X & SP-Exp2 & 240~s & 9.699 & 154.8~s & early \\
X & SP-Exp3 & 120~s & 2.142 & no; peak 204.9\textdegree C & late \\
X & SP-Exp3 & 120~s & 3.103 & no; peak 236.8\textdegree C & late \\
X & SP-Exp3 & 120~s & 4.064 & 193.8~s & early \\
X & SP-Exp3 & 240~s & 2.528 & no; peak 217.5\textdegree C & late \\
X & SP-Exp3 & 240~s & 3.662 & 242.1~s & late \\
X & SP-Exp3 & 240~s & 4.229 & 221.5~s & late \\
X & SP-Exp3 & 240~s & 4.370 & 217.2~s & early \\
X & SP-Exp3 & 240~s & 4.512 & 213.1~s & early \\
X & SP-Exp3 & 240~s & 4.795 & 205.7~s & early \\
X & ST-Exp4 & 120~s & 2.933 & no; peak 197.2\textdegree C & late \\
X & ST-Exp4 & 120~s & 3.926 & no; peak 209.1\textdegree C & late \\
X & ST-Exp4 & 120~s & 4.422 & no; peak 224.2\textdegree C & late \\
X & ST-Exp4 & 120~s & 4.547 & no; peak 238.3\textdegree C & late \\
X & ST-Exp4 & 120~s & 4.671 & 292.6~s & early \\
X & ST-Exp4 & 120~s & 4.919 & 231.3~s & early \\
X & ST-Exp4 & 240~s & 4.103 & no; peak 231.2\textdegree C & late \\
X & ST-Exp4 & 240~s & 4.292 & 316.5~s & early \\
X & ST-Exp4 & 240~s & 4.481 & 278.5~s & early \\
\bottomrule
\end{longtable}
}
}{}

\section{Results quoted without another table}
\label{sec:s-quoted}

Table~\ref{tab:s-quoted} gathers every result that the main text quotes from a run or a computation that no other table or figure shows, and every result of a run that the text of this supplement quotes and no other table gives, each with the runs or the code it comes from, and the inputs that the earlier study set without a primary source.

\IfFileExists{tables/s17_quoted.tex}{{\scriptsize\setlength{\tabcolsep}{4pt}
\begin{longtable}{>{\raggedright\arraybackslash}p{0.07\textwidth}>{\raggedright\arraybackslash}p{0.41\textwidth}>{\raggedright\arraybackslash}p{0.17\textwidth}>{\raggedright\arraybackslash}p{0.25\textwidth}}
\caption{Every result that the main text (by section) quotes from a run or a computation that no other table or figure shows, every result of a run that the text of this supplement (S2, S3) quotes and no other table gives, and the inputs that the earlier study set without a primary source (Table~2 of the main text, $^{a}$), each with the runs or the code it comes from. Every value here is read from its file, by the numbers check's own entry for it where a text quotes it.}\label{tab:s-quoted}\\
\toprule
Where & Result & Value & From \\
\midrule
\endfirsthead
\multicolumn{4}{l}{\textit{Continued}} \\
\toprule
Where & Result & Value & From \\
\midrule
\endhead
2.1 & The median ratio of ejected to body energy over the Databank's ten nail tests of this cell, by which each test's ejected energy is taken from its fitted body energy & 3.105 & the Databank's rows; \nolinkurl{results/stage8_ejected.csv}, ratio\_central \\
2.3, S2 & The posterior-mean kinetics in the emulation of the kinetic source's accelerating-rate calorimetry: the exotherm onset, and the temperature at which the self-heating reaches 1 K\,min$^{-1}$ & 179.3 \textdegree C; 199.8 \textdegree C & the emulation; \nolinkurl{results/note_numbers.json} \\
2.3, S2 & The 200 kinetic draws in the same emulation: how many reach that rate within 4,000 min (the source's main cluster), their exotherm onsets (central 95\,\%) and mean, and the band of the check as amended, their mean and two standard deviations & 164 of 200; 146.3 to 191.4 \textdegree C; 178.8 \textdegree C; 161.9 to 195.7 \textdegree C & the emulation; \nolinkurl{results/stage3_v4.json}, \nolinkurl{results/stage3_v4_amended.json} \\
3.1 & With body energies up to 33.2 kJ and no delivered heat, over the 5,248 runs at every contact, tab conductance, release time and connection and under the 200 kinetic draws: the highest peak of the resolved neighbour's hottest point, its margin to 255.7 \textdegree C, and the lumped neighbour's highest peak & 226.7 \textdegree C; 29.0 K; 180.6 \textdegree C & step 5's runs; \nolinkurl{results/v12_numbers.json}, all\_runs \\
3.1 & The eight tests with the largest fitted triggers held to 33.2 kJ, no delivered heat: the hottest point's peak at the central and at the high contact (two connections each) and the lumped neighbour's highest peak & 191.9 and 192.2 \textdegree C; 226.4 and 226.7 \textdegree C; 177.4 \textdegree C & step 12's runs; \nolinkurl{results/stage12_paper_numbers.json}, part\_a.central and part\_a.high \\
3.1 & With the fitted triggers and no delivered heat, the latest time by which conduction carries the hottest point over 255.7 \textdegree C in the eight tests with the largest triggers & 21.7 s & step 12's runs; \nolinkurl{results/stage12_paper_numbers.json}, part\_a.central \\
3.1 & With the fitted triggers and no delivered heat at the low contact, the highest peak of the hottest point over the sixteen tests & 162.6 \textdegree C & step 7's runs; \nolinkurl{results/stage7_crossings.csv} \\
3.2 & With no delivered heat at the central contact, how far the far face's peak exceeds the measured maximum in the four tests that did not propagate and need none & 13.8 to 56.3 K & step 9's runs; smallest and largest of \nolinkurl{results/stage9_paper_numbers.json}, part\_a.A3 \\
3.3 & PT-Exp1 with 11.32 kJ over 21.4 s, the upper end of its three-dimensional threshold's bracket: the lumped neighbour's peak & 204.2 \textdegree C & step 16's run; \nolinkurl{results/stage16_runs.csv} \\
3.3 & ST-Exp3 with 10.80 kJ over 21.4 s: the resolved neighbour's runaway & 61.5 s & step 15's run; \nolinkurl{results/stage15_paper_numbers.json}, B1 \\
3.4 & PT-Exp1 with a 360 s pulse at the upper end of the local threshold's bracket: the hottest point's crossing, on the cell's axis, and the mean then & 316.5 s; 240.1 \textdegree C & step 13's run; \nolinkurl{results/stage13_skin.csv} \\
3.4 & The lumped neighbour's ignition temperature with the trigger at the neighbour's own temperature, and there, with every reactant at its initial amount, the share of the low-temperature cathode reaction in the reaction heat under the posterior mean & 218.5 \textdegree C; 73.3 \% & step 13's computation; \nolinkurl{results/stage13_tign.csv} (trigger own) and \nolinkurl{results/stage13_paper_numbers.json} \\
3.4 & With no delivered heat in PT-Exp4, PP-Exp3 and ST-Exp2 at the central contact, the heat the lumped neighbour has taken in through the contact and wedge by 250 s beyond the resolved one's & 2.1 to 2.3 kJ & step 15's accounts of archived runs; \nolinkurl{results/stage15_accounts.json} \\
3.4, S2 & PT-Exp4 with 5.04 kJ over 21.4 s: the lumped neighbour's runaway under the posterior mean and under the kinetic draws ranked 5 and 50 & 215.6 s; 191.0 and 196.5 s & step 15's runs; \nolinkurl{results/stage15_readings.json} and \nolinkurl{results/stage15_paper_numbers.json}, B4 \\
3.4 & PT-Exp4 with 5.04 kJ over 21.4 s: the resolved neighbour's mean peak under the four kinetic draws ranked 5, 50, 151 and 196, and with every rate law shifted by 17.5 K & 191.2 to 192.4 \textdegree C; 221.1 \textdegree C & step 15's runs; \nolinkurl{results/stage15_paper_numbers.json}, B4 \\
3.4 & PT-Exp4 with 5.04 kJ over 21.4 s and every rate law shifted by 20 K: the resolved neighbour's runaway & 270.3 s & step 15's run; \nolinkurl{results/stage15_paper_numbers.json}, B4 \\
3.4 & The lumped threshold with the trigger's body energy shifted by 2 kJ down and up (Fig. 4b of the main text), and its fall for each kJ added, the difference of the two brackets' midpoints over 4 kJ & PP-Exp3 (3.83, 3.89] and (2.02, 2.06] kJ, 456 J; ST-Exp2 (3.04, 3.08] and (1.31, 1.33] kJ, 435 J; PT-Exp3 (8.93, 9.07] and (7.37, 7.48] kJ, 393 J & step 13's runs; \nolinkurl{results/stage13_scorecard.json}, P10 cases \\
3.5 & PT-Exp4, the 5.04 kJ pulse followed by 10 W and by 12.5 W: the resolved neighbour's runaway, and with 12.5 W the far face's reading where its jump begins & 369.3 and 255.9 s; 233.6 \textdegree C & step 15's runs; \nolinkurl{results/stage15_paper_numbers.json}, B3b \\
3.5 & PT-Exp4 with 7.42 kJ over 21.4 s: how far the far face's jump begins above the measured 197.3 \textdegree C, and above the record's own jump read by the same rule & 36.8 K; 47.6 K & step 15's run; \nolinkurl{results/stage15_scorecard.json} and \nolinkurl{results/stage15_paper_numbers.json} \\
3.5 & PT-Exp4 with 6.2 kJ at constant power over 240 s, the run of step 10 carried on to 900 s: the resolved neighbour's mean peak & 241.4 \textdegree C & step 15's run; \nolinkurl{results/stage15_paper_numbers.json}, B2 \\
3.5 & PT-Exp4: the root mean square of the far face's difference from the record over 0 to 271.5 s with 4.66 and 5.04 kJ over 21.4 s, and over 0 to 181.9 s with 7.42 kJ & 16.3 and 16.8 K; 51.4 K & steps 10 and 15; \nolinkurl{results/stage10_paper_numbers.json}, \nolinkurl{results/stage15_paper_numbers.json} \\
3.7 & Wall time of a run on one two-core machine: the reactive cross-section, and the three-dimensional model in the same searches & 17 to 71 s; 246 to 512 s & \nolinkurl{results/stage15_xs/*.json} and \nolinkurl{results/stage15_runs/*.json}, wall\_s \\
App. B & The reduced models' constants: the roll's effusivity across its layers, the can's heat capacity per unit of the cell's outer area, the skin model's factor $\Phi$, and the whole heat of the low-temperature cathode reaction at the posterior mean & 1,743 J\,m$^{-2}$\,K$^{-1}$\,s$^{-1/2}$; 841.7 J\,m$^{-2}$\,K$^{-1}$; 1.2284; 2,667 J & step 13's computation; \nolinkurl{results/stage13_readings.json}, \nolinkurl{results/stage13_checks.json}, \nolinkurl{results/stage13_tign.csv} \\
S2 & Step 3, the earlier geometry's case (33.2 kJ over 30 s, a 500 mW\,K$^{-1}$ busbar): the hot spot's onset over the 200 kinetic draws, under the four draws step 11 took, under the posterior mean, and its standard deviation over the kinetic and over the 30 thermal draws & 6.01 to 17.13 s; 6.25, 10.23, 16.04, 17.02 s; 11.48 s; 2.83 and 0.076 s & step 3's runs; \nolinkurl{results/stage3a_onsets.csv}, stage3a\_metrics.json, stage11\_draws.csv, stage3bc\_metrics.json \\
S3 & Step 5's central case (the central contact and tab, the posterior-mean kinetics), over its two releases and two connections: the hottest point's peak and the lumped neighbour's & 191.5 to 192.3 \textdegree C; 159.4 to 162.1 \textdegree C & step 5's runs; \nolinkurl{results/v12_numbers.json}, central\_case.span \\
S3 & Step 9's runs with the pulse at or above the lower end of each propagated test's local threshold: the latest crossing of 255.7 \textdegree C by the hottest point, after the trigger's start & 13.8 to 21.1 s & step 9's runs; \nolinkurl{results/stage9_scorecard.json}, P1 rows \\
S3 & Step 6's runs with 33.2 kJ over 21.39 s under the posterior mean: the hottest point's peak at the central contact and tabs in series and in parallel, and at the high contact with the low and the high tabs & 192.16 and 191.90 \textdegree C; 226.14 to 226.75 \textdegree C & step 6's runs; \nolinkurl{results/stage6_runs/*E33.22.json}, resolved.peak\_hottest\_C \\
Table 2 & Tab conductance per tab, set by the earlier study without a primary source and held here as an assumption & 6.8 mW\,K$^{-1}$ (5 and 10 as brackets) & \nolinkurl{src/v12_model.py} line 37 (G\_TAB = 0.005, 0.0068, 0.010 \nolinkurl{W/K}) \\
Table 2 & Can conductivity, set by the earlier study without a primary source and held here as an assumption & 15 W\,m$^{-1}$\,K$^{-1}$ & \nolinkurl{external/m50-two-cell-pilot/src/lumped_two_cell.py} line 57 (K\_CAN = 15.0) \\
Table 2 & Can mass and specific heat, set by the earlier study without a primary source and held here as an assumption & 9 g; 500 J\,kg$^{-1}$\,K$^{-1}$ & \nolinkurl{external/m50-two-cell-pilot/src/lumped_two_cell.py} lines 54 to 55 (M\_CAN, CP\_CAN) \\
Table 2 & Air in the wedge, set by the earlier study without a primary source and held here as an assumption & 0.03 W\,m$^{-1}$\,K$^{-1}$ & \nolinkurl{external/m50-two-cell-pilot/src/lumped_two_cell.py} line 69 (K\_AIR = 0.03) \\
Table 2 & Contact sector, either side of the line of contact, set by the earlier study without a primary source and held here as an assumption & 45\textdegree & \nolinkurl{src/v12_mesh.py} line 22 (SECTOR = pi / 4, "half-angle" \\
Table 2 & Room convection, set by the earlier study without a primary source and held here as an assumption & 10 W\,m$^{-2}$\,K$^{-1}$ & \nolinkurl{external/m50-two-cell-pilot/src/lumped_two_cell.py} lines 63 and 74 (H\_CONV, A\_CONV) \\
Table 2 & Cell mass, set by the earlier study without a primary source and held here as an assumption & 65 g & \nolinkurl{external/m50-two-cell-pilot/src/lumped_two_cell.py} line 47 (M\_CELL = 0.065) \\
\bottomrule
\end{longtable}
}
}{}

\bibliographystyle{elsarticle-num}
\bibliography{references}

\end{document}
